% MSc dissertation example file, February 2022
%
% Leave one of the documentclass lines uncommented to match your degree.
% You may remove the logo option if it causes problems.
% Do not change any other options.
% \documentclass[logo,msc,adi]{infthesis}     % Adv Design Inf
% \documentclass[logo,msc,ai]{infthesis}      % AI
% \documentclass[logo,msc,cogsci]{infthesis}  % Cognitive Sci
% \documentclass[logo,msc,cs]{infthesis}      % Computer Sci
% \documentclass[logo,msc,cyber]{infthesis}   % Cyber Sec
% \documentclass[logo,msc,datasci]{infthesis} % Data Sci
% \documentclass[logo,msc,di]{infthesis}      % Design Inf
% \documentclass[logo,msc,dsti]{infthesis}    % Data Sci TI
% \documentclass[logo,msc,inf]{infthesis}     % Informatics
\documentclass[logo,msc,ai]{infthesis}           % degree unspecified, do not change except to add your degree
%%%%%%%%%%%%%%%%%%%%%%%%
% Understand any problems and seek approval before assuming it's ok to remove ugcheck.
% \usepackage{msccheck}
\usepackage{msccheck} 

% Include any packages you need below, but don't include any that change the page
% layout or style of the dissertation. By including the ugcheck package above,
% you should catch most accidental changes of page layout though.

\usepackage{microtype} % recommended, but you can remove if it causes problems
\usepackage{xcolor} %remove later this is used for helping edit and draft
\usepackage{caption}
\usepackage{subcaption} %subcaptions for plots and figures next to each other. 
\usepackage{amsmath}  %Check we are allowed
\usepackage{amssymb} %Check we are allowed
\usepackage{amsfonts} %Check we are allowed
\usepackage{booktabs} %Check we are allowed
\usepackage{graphicx} %Check we are allowed
\usepackage{tabularx}
\usepackage{float}
\usepackage{underscore} %Check we are allowed
\usepackage{longtable} % lets a table break across pages instead of overflowing one page


%=======================================================================================
\begin{document}
\begin{preliminary}

\title{Physics-informed prediction and spatial attribution of forest top-height growth using repeated LiDAR surveys}
%TODO: DRAFT DIFFERENT VERSIONS OF THE TITLE AND PICK AGAIN AT THE END

\author{Shreeya Gorasia}

\date{\today}
%=======================================================================================

\abstract{
Forestry models rely on standard growth curves to show how stands develop on average, but miss why individual trees stray from that path. This dissertation combines repeated airborne LiDAR, forest records, and environmental data to track Sitka spruce top-height growth across Aberfoyle Forest, Scotland, spanning 15 survey years.

Linear models, mixed-effects models, and XGBoost identified factors linked to each plot's average gap from a baseline Chapman--Richards curve. A spatially varying neural model (GNNWR) tested whether these links shift across the landscape by targeting the gap between a plot's fitted height ceiling and its yield-class reference. Standard machine learning, neural networks, and physics-guided models were then compared for direct top-height prediction using compartment-level spatial cross-validation to test whether adding environmental features helps.

Canopy cover was the strongest driver of growth gaps. Management added minor signals, while broad environmental layers added little beyond local wind shelter and exposure effects. Although GNNWR slightly beat global models, canopy cover drove 80--90\% of its signal, showing little extra spatial pattern from the environment. XGBoost gave the most accurate top-height predictions. PINN-$k$ had an RMSE roughly 0.4\,m higher, but it eliminated impossible predictions (5--70\,m) while producing meaningful, plot-level parameters.

Local stand structure gave the clearest signals, while spatial models added little value. Physics guidance provided interpretable parameters and biological safety rather than boosting raw accuracy.
}

%TODO: writing: Comeback to this section when the entire report is done 
%=======================================================================================

\maketitle

\newenvironment{ethics}
   {\begin{frontenv}{Research Ethics Approval}{\LARGE}}
   {\end{frontenv}\newpage}

\begin{ethics}
This project was planned in accordance with the Informatics Research
Ethics policy. It did not involve any aspects that required approval
from the Informatics Research Ethics committee.

I declare that this thesis was composed by myself, that the work contained herein is my own except where explicitly stated otherwise in the text, and that this work has not been submitted for any other degree or professional qualification except as specified. AI tools were used throughout this project to help plan work, research papers, brainstorm ideas, edit code for plotting, and submit cluster jobs. I always checked this assistance myself, tested it, and critically evaluated its output before relying on it, so the ideas and writing in this thesis are my own. AI assistance for writing was limited to feedback and suggestions, not direct edits to the submitted text.
\end{ethics}

%TODO: writing: Comeback to this section when the entire report is done 
%=======================================================================================

\begin{acknowledgements}
My sincere thanks go first to my supervisor, Dr Michael Herrmann, for his guidance, thoughtful feedback, and continued support throughout this dissertation. His advice helped shape the direction of the project and resolve the challenges that arose along the way.

I am also very grateful to Dr Juan Suárez and Miguel Ibáñez at Forest Research for generously sharing their time and expertise. Their help in understanding the forest data, clarifying its wider forestry context, and patiently answering my many questions was invaluable to this work.

Special thanks go to the friends I have made in Edinburgh. Studying alongside all of you made this degree more enjoyable and memorable. Your encouragement, patience, and willingness to listen to countless ideas and dissertation updates meant a great deal to me. I could not have completed this project without your support. Above all, you made my time in Edinburgh truly special, both inside and outside the library.

Finally, my deepest thanks go to my family. Your belief in me has carried me through every stage of my education. Your encouragement during difficult moments, and your reminders of how far I had already come, meant more than I can express. I am deeply grateful for the sacrifices you have made and for the love and support that made this Master's degree possible.

\end{acknowledgements}

%TODO: writing: Comeback to this section when the entire report is done 
%=======================================================================================

\tableofcontents
\end{preliminary}

%TODO: check: Comeback to this when entire report is done and check it and make it fits in one page
%=======================================================================================
\chapter{Introduction}  \label{ch:intro}

\section{Motivation and study context}
Forest Digital Twins are a wider long-term goal for the forestry sector: digital models of forest systems that help monitor, forecast, and  management planning \cite{baskent2026Review_C1_SS, buonocore2022Proposal_C1_SS}. They aim to replace static, periodic checks with adaptive, data-driven planning by combining repeated field observations, simulation, and machine learning(ML), and artificial intelligence(AI) \cite{baskent2026Review_C1_SS, buonocore2022Proposal_C1_SS}. This is vital as managers face climate risks—like heavy storms, high winds, and severe summer heatwaves that spark wildfires and shift growth. At the same time, forestry is moving towards continuous cover management, which demands more flexible stand-level models \cite{ForestResearch2009Evidence_C1_SS}.

Sitka spruce (\textit{Picea sitchensis}) is Britain's primary commercial species, making up approximately 50\% of managed forests \cite{ForestResearch2026Sitka_C1_SS} and contributing roughly \pounds1.1 billion annually to Scotland's economy \cite{ScottishForestry2026Future_C1_SS, mason2011Sitka_C1_SS}. The study focuses on Aberfoyle in Trossachs, Scotland, a Sitka spruce-dominated region with 21 years of repeated LiDAR coverage. Featuring varied topography with exposed, windy ridges to sheltered valleys, providing variation that causes different growth trajectories in different areas of the forest.

% Height, LiDAR and shared growth relationship
Top height is a key measure of stand development, commonly used in forestry models to estimate downstream metrics like timber volume and site productivity \cite{ForestResearch2026Sitka_C1_SS}, and now can be accurately captured at scale using airborne LiDAR scans \cite{suarezminguez2008Practical_C1_SS}. For most tree species, including Sitka spruce, height development over time is modelled by the Chapman-Richards equation \cite{pommerening2015Methods_C2_CR}, a sigmoid growth curve, with its parameters fixed to represent a stand's maximum height potential and its speed of growth. A fixed parameter curve only captures the average population trend, but local trajectories can deviate from this baseline, due to factors like terrain, climate, soil, wind exposure, and active management all critically constrain or enhance a stand's actual development \cite{socha2023Higher_C1_SS, cameron2015Building_C1_SS}.  This raises two key questions: whether local departures from this shared curve are associated with identifiable environmental and management conditions, and whether independent attribution methods agree on which conditions matter.

Foresters need to know if a stand's assigned yield class still holds, or which local factors push growth off track. This dissertation tests whether letting these environment--growth relationships vary across space captures local growth patterns that fixed regional curves miss. Flexible machine learning models can also learn these patterns directly from data without fixed curve assumptions; we test deep neural networks (DNNs) and gradient-boosted trees (XGBoost) as a supplementary check on top-height predictions.

\section{Research gap and contributions}
While repeated airborne laser scanning has been applied to model forest height growth and productivity \cite{janiec2023Development_C1_SS}, limited evidence exists on whether local environmental conditions are associated with departures from Sitka spruce's shared height-growth curve, and whether different attribution methods identify consistent conditions. Many current remote-sensing models do not account for spatial dependence between nearby-identical trees, leading to data leakage and inflated performance estimates \cite{Karasiak2022Spatial_C5_EVAL, ploton2020Spatial_C5_EVAL, lynch2025Digital_C4_NN}. Relying on black-box neural networks without explainable AI limits their use for stakeholders who require interpretable insights to support decision-making.

This dissertation covers one part of the Digital Twin idea: using attribution methods to explain local growth deviations, and checking if that explanation can help predict top height. It tests whether growth-environment links change by location, using a Geographically Neural Network-Weighted Regression (GNNWR) model, and compares standard and physics-informed neural networks (PINNs) on prediction itself, to see which one wins and why.

\section{Objectives and research questions}
Answering the three research questions required combining repeated forest surveys, airborne LiDAR, and environmental rasters (terrain, wind, climate, and soil) into a single dataset. Together, they form a clear workflow for forest management: Q1 flags stands straying from the standard regional growth curve and finds the environmental drivers behind it; Q2 tests whether a stand's assigned yield-class ceiling ($\Delta y_{\max}$) needs local adjustment; and Q3 tests whether adding these site features and physics rules improves top-height predictions in the field.

\begin{table}[!htb]
\small
\setlength{\tabcolsep}{5pt}
\renewcommand{\arraystretch}{1.08}
\begin{tabularx}{\linewidth}{@{} >{\raggedright}p{0.32\linewidth} >{\raggedright\arraybackslash}X @{}}
\noalign{\hrule height 1pt}
\textbf{Research Question \& Target} & \textbf{Objectives \& Operational Purpose} \\
\noalign{\hrule height 0.6pt}
\textbf{Q1: Baseline Curve Attribution} \newline
\textit{Target:} Multi-year mean residual ($\bar{r}_i$) from the shared Chapman--Richards curve & 
\textbf{(1.1)} Link persistent plot-level growth deviations to terrain, soil, wind, and management. \newline
\textbf{(1.2)} Test if drivers agree across Elastic Net, LMM, and XGBoost. \newline
\textbf{Purpose:} Flags off-track stands and tests attribution consistency across model types. \\
\hline
\textbf{Q2: Spatial Variation in Growth} \newline
\textit{Target:} Yield-class asymptote deviation ($\Delta y_{\max}^{(i)}$) & 
\textbf{(2.1)} Test spatially varying models (GNNWR) under spatial CV to capture local patterns. \newline
\textbf{(2.2)} Separate genuine local environmental trends from noise and fitting artefacts. \newline
\textbf{Purpose:} Diagnose if a stand's standard yield class needs revision based on local site effects. \\
\hline
\textbf{Q3: Prediction \& Physics Guidance} \newline
\textit{Target:} Direct LiDAR top height ($H_{it}$) &
\textbf{(3.1)} Benchmark empirical models against two physics-informed neural networks: PINN (height only) and PINN-$k$ (height and growth rate). \newline
\textbf{(3.2)} Test if environmental features and physics constraints improve predictions or parameter interpretability, and compare the two PINN versions. \newline
\textbf{Purpose:} Finds the best operational model for future height prediction and tests if PINN-$k$'s extra parameter justifies its small drop in accuracy. \\
\noalign{\hrule height 1pt}
\end{tabularx}
\caption{\small{Overview of research questions, objectives, and wider management purpose.}}
\label{tab:aims_objectives}
\end{table}


\paragraph*{Main findings}
Environmental and management conditions can explain some of the growth differences within Aberfoyle. Elastic Net, LMM, and XGBoost all point to similar features (Q1). GNNWR finds little that these three miss; most local patterns are just noise, or flaws in how the target was built (Q2). XGBoost is the most accurate model for predicting top height, outperforming the standard DNN and both PINN variants; the two-parameter PINN loses a little accuracy but never predicts an impossible height (Q3). Physics guidance and local spatial modelling therefore do not improve accuracy, but yield safer, more interpretable models.


%=======================================================================================
%=======================================================================================

\chapter{Related Work and Background} \label{ch:background}

\section{Shared height-growth models}
Top height tracks upper-canopy development and serves as a standard metric of forest growth and site productivity. Multi-temporal airborne LiDAR surveys now allow top height to be estimated consistently across broad landscapes, though uncertainty persists in young stands, open canopies, and edge boundaries \cite{ForestResearch2026Sitka_C1_SS, suarezminguez2008Practical_C1_SS}

The Chapman-Richards(CR) model describes height trajectories using a flexible sigmoid curve parametrised by an asymptotic ceiling ($y_{\max}$), an intrinsic growth rate ($k$), and a shape exponent ($p$). General Yield Class (GYC) groups stands by their expected timber production rate, with each class mapped to a standard height growth curve. While a single population curve offers a useful baseline, it cannot capture varying local site conditions and management histories \cite{chapman1961Statistical_C2_CR, pommerening2015Methods_C2_CR, mason2011Sitka_C1_SS}

\begin{figure}[!htbp]
    \centering
    \includegraphics[width=0.65\textwidth]
    {figures/fig_background/CR-curves-yeildclass.png}
    \caption{Chapman-Richards growth curves coloured by General Yield Class (GYC), work done by Forest Research}
    \label{fig:chapman-richards-illustrative}
\end{figure}

\section{Environmental controls on growth}
Sitka spruce productivity responds strongly to elevation, wind exposure, soil moisture, and soil fertility: in upland northern Britain, General Yield Class falls by roughly $3.2\text{-}4.0\,\mathrm{m^3\,ha^{-1}\,yr^{-1}}$ per 100\,m rise in elevation, from lower temperatures and more wind exposure \cite{worrell1990Productivity1_C1_SS}, and Irish plantations show the same pattern, peak growth on moist, fertile sites, with sharp declines on drought-prone or waterlogged soils \cite{farrelly2011Sitka_C1_SS}.

Landscape-scale studies combine forest inventory with spatial layers, from early UK work pairing sub-compartment records with terrain and soil maps via regression \cite{bateman1998Using_C1_SS} to more recent studies adding LiDAR, climate, and soil-type rasters. Since predictor importance shifts with spatial resolution and extent \cite{fricker2019More_C3_LID}, and coarse soil layers often match local conditions poorly \cite{miller2024Assessing_C2_CR}, differences between studies likely reflect study design rather than one universal set of drivers.

\section{Global growth attribution}
Global models apply one environment-growth rule across the entire forest rather than letting it change by location. Mixed-effects models are common in forestry for separating broad environmental trends from stand-level differences \cite{harrison2018Brief_C5_EVAL, socha2016Assessment_C2_CR}. However, their random effects merely soak up leftover plot variance rather than explaining its ecological causes. Elastic Net handles correlated variables well by combining penalty terms ($L_1$ and $L_2$), but it still assumes relationships are strictly linear and additive \cite{zou2005Regularization_C4_NN}. Because each method relies on different mathematical assumptions, comparing these model families is a common way to test whether an association is robust, or an artefact of one model's own structure.

Boosted-tree models have been used to map forest productivity from LiDAR and environmental data, while XGBoost-SHAP has been applied to climatic and environmental drivers of forest condition \cite{puliti2023Tree_C2_CR, chen2016XGBoost_C4_NN}. When paired with tree models, SHAP has been used to identify non-linear thresholds and interactions learned by fitted tree models \cite{lundberg2017Unified_C7_XAI}. However, SHAP values reflect feature attribution within the model rather than direct biological causality, and attribution can be redistributed among correlated predictors \cite{aas2021Explaining_C7_XAI}. Limited work compares linear, mixed-effects, and tree-based attribution on the same persistent growth-departure target.

\section{Spatially varying attribution}
Spatial autocorrelation occurs when nearby measurements are similar to one another. By contrast, spatial non-stationarity means that the relationship between the environment and tree-growth changes from one geographic area to another. 

When a model leaves residual spatial autocorrelation, neighbouring areas are systematically over- or under-predicted. This may reflect missing spatial predictors, an oversimplified functional form, sampling effects, or genuine spatial non-stationarity \cite{gaspard2019Residual_C6_GWR}.

Geographically Weighted Regression (GWR) represents non-stationarity by fitting local regressions in which nearby observations receive greater weight. Its local coefficients can be mapped to show how estimated relationships vary across space \cite{brunsdon1996Geographically_C6_GWR}. However, local multicollinearity can produce unstable or strongly correlated coefficients, making these maps difficult to interpret \cite{wheeler2005Multicollinearity_C6_GWR}. Multiscale GWR (MGWR) relaxes GWR's single-scale assumption by allowing each predictor to operate over a different spatial bandwidth \cite{fotheringham2017Multiscale_C6_GWR}.

Geographically Neural Network-Weighted Regression (GNNWR) replaces predefined spatial kernels with a neural network that learns the relationship between distance and local weighting \cite{du2020Geographically_C6_GWR}. This addresses limitations of classical GWR, whose results depend on the selected bandwidth and kernel form and whose simple weighting functions may not represent complex spatial processes. In forestry, GNNWR improved prediction and reduced residual spatial autocorrelation when relating down wood to terrain and stand characteristics \cite{sun2022Local_C6_GWR}. However, evidence came from one forest reserve, one disturbance outcome, and five retained predictors, while the study gave limited attention to GNNWR's own interpretability and generalisability.


\section{Physics-guided prediction}
Physics-informed neural networks use governing equations as soft constraints during training \cite{raissi2019Physics_C4_NN}. Where exact physical laws are unavailable, process-guided models instead embed approximate biological knowledge. Forest applications have represented linked processes such as growth, competition, regeneration and mortality \cite{pichler2025Inferring_C4_NN}. However, rigid process assumptions can bias estimates of other processes, while sparse observations may leave alternative mechanisms difficult to distinguish. Even flexible hybrid models can learn biologically implausible responses in poorly sampled conditions \cite{pichler2025Inferring_C4_NN}.

Training also requires empirical and process losses to be balanced. Numerical stiffness can produce unequal gradients, causing one objective to dominate optimisation \cite{wang2021Understanding_C4_NN}. Chapman-Richards guidance therefore encourages an established growth shape but does not guarantee either biological validity or improved prediction. Although process-guided networks are widely studied, integrating Chapman-Richards equations into neural networks remains rare outside specialised variable-age growth models but it has not been tested under long-term LiDAR time series in Aberfoyle.

\section{Spatial model evaluation}
Residual Moran's $I$ is used in ecology to spot spatial patterns that average error scores can miss. It gives a simple test of whether nearby plots are consistently over- or under-predicted, indicating if a model has missed key spatial structure or yields misleading conclusions \cite{moran1950Notes_c5_eval,kim2021Predicting_C5_EVAL}. In forestry, it helps compare global, GWR, and neural spatial models, where local neural models produced both lower errors and less residual clustering \cite{sun2022Local_C4_NN}. However, Moran's $I$ relies on the chosen distance settings and compresses the whole map into a single number, which can hide distinct local patches. A lower score shows less error clustering under those settings, but it cannot explain why or prove that the model found a true biological relationship.

\section{Research gap and study position}
Existing studies show that the environment strongly shapes forest growth, but they typically treat growth curves, driver attribution, spatial variation, and prediction as separate problems. No previous study directly compares linear, mixed-effects, non-linear, and spatial attribution models using multi-temporal LiDAR trajectories on the same dataset with spatially independent validation. As a result, it remains unclear whether identified drivers agree across model types, whether their effects change across the landscape, and whether adding this environmental information actually improves top-height predictions.

This work addresses these gaps through three contributions:
\begin{enumerate}
\setlength\itemsep{0pt}
\item \textbf{Global attribution (Q1):}
Compares whether environmental associations with persistent curve departure remain consistent across Elastic Net, mixed-effects and XGBoost models.
\item \textbf{Spatial attribution (Q2):}
Tests whether associations with fitted asymptotic departure vary geographically using GNNWR and spatially blocked evaluation.
\item \textbf{Prediction and process guidance (Q3):}
Compares standard and neural top-height models, then tests whether environmental inputs or Chapman--Richards guidance improve accuracy or parameter interpretation.
\end{enumerate}
%=======================================================================================
%=======================================================================================

%=======================================================================================
%=======================================================================================
\chapter{Data}
This chapter describes the dataset used for model training and evaluation, outlining key feature engineering steps and data limitations. Table~\ref{tab:data-fields} summarises all predictor and target variables alongside their sources and temporal coverage. Feature selection rationale, data partitioning, and specific modelling choices are detailed in Chapter~\ref{ch:methodology}.

\section{Study area and cohort coverage}
The dataset covers approximately $16\text{~km} \times 12\text{~km}$ of Aberfoyle and contains 795,645 plot--survey records from 230,359 unique plot locations across six survey years (2002--2023). The records contain airborne laser scanning (ALS)-derived variables, including canopy cover and height percentiles, alongside stand and management fields such as planting year and known felling history (Table~\ref{tab:data-fields}). Figure~\ref{fig:data-study-area} shows the study-area topography, the spatial coverage of both cohorts, and an example of standard and boundary-clipped plots.

\begin{figure}[H]
\includegraphics[
    width=1.05\textwidth,
    height=0.48\textheight,
    keepaspectratio
]{figures/fig_data/aberfoyle_area.png}
\caption{\small{Left:Aberfoyle study area. Middle: cohort coverage, Right: plot-level detail}}
\label{fig:data-study-area}
\end{figure}

Because LiDAR coverage increase by $37\text{ km}^2$ between 2006 and 2008, two nested cohorts were defined (Table~\ref{tab:data-cleaning-funnel}): a broad-coverage \emph{4-survey} cohort (2008, 2012, 2021, 2023; 71,766 plots across 296 compartments) and a longer-trajectory \emph{6-survey} cohort (2002--2023; 13,897 plots across 48 compartments).

To ensure data quality, sequential filters were applied to both cohorts (Table~\ref{tab:data-cleaning-funnel}):
\begin{itemize}
    \item \textbf{Complete presence:} Plots must be present across all survey years in the cohort.
    \item \textbf{Species purity:} Pure Sitka spruce stands only.
    \item \textbf{Plausible records:} Planting year between 1850 and the survey year, stand age $0\text{--}200$ years, and LiDAR top height $0\text{--}60\text{ m}$ (ecological maximum \cite{mason2011Sitka_C1_SS}).
    \item \textbf{Maturity cut-off:} Stand age $\ge 30$ years by 2023 to exclude unreliable young-stand height measurements \cite{suarezminguez2008Practical_C1_SS}.
\end{itemize}

\begin{table}[!htbp]
\centering
\scriptsize
\begin{tabular}{lrrrr}
\toprule
& \multicolumn{2}{c}{4-survey} & \multicolumn{2}{c}{6-survey} \\
\cmidrule(lr){2-3}\cmidrule(lr){4-5}
Stage & Plots & Removed & Plots & Removed \\
\midrule
0. Raw dataset
& 230,359 &, & 230,359 &, \\
1. Cohort years present
& 230,345 & 0.01\% & 230,359 & 0.00\% \\
2. Complete survey coverage
& 112,782 & 51.04\% & 18,781 & 91.85\% \\
3. Sitka spruce only
& 72,395 & 35.81\% & 13,897 & 26.01\% \\
4. Valid planting year
& 71,766 & 0.87\% & 13,897 & 0.00\% \\
5. Plausible age (0-200 years)
& 71,766 & 0.00\% & 13,897 & 0.00\% \\
6. Plausible height (0-60 m)
& 71,766 & 0.00\% & 13,897 & 0.00\% \\
7. Mature stands (age $\geq 30$ in 2023)
& \textbf{58,112} & 19.03\%
& \textbf{13,769} & 0.92\% \\
\bottomrule
\end{tabular}
\caption{\small{Plots left after each filter. Percentages show the share removed from the previous row, rather than from the raw dataset. Figure~\ref{fig:data-study-area} shows the final spatial coverage.}}
\label{tab:data-cleaning-funnel}
\end{table}

%------------------------------------------------
\section{Response variables and targets}
All targets in this dissertation derive from the plot-level 95th percentile ALS return height (\texttt{elev\_percentile\_95th}), forming three distinct response variables. The multi-year mean residual from a shared growth curve (Q1, $\bar{r}_i$), the gap between a plot's fitted ALS height ceiling and its assigned yield-class ceiling (Q2, $\Delta y_{\max}^{(i)}$), and direct LiDAR top height (Q3, $H_{it}$). Forest Research records supply planting year, management history, and yield class (Table~\ref{tab:data-fields}), with any variables calculated directly from measured height removed from model inputs to prevent target data leakage. A code-level audit confirmed that the pipeline contains no target circularity: environmental features (terrain, climate, wind, and soil) were extracted independently from external spatial layers and joined only at the modelling stage, never entering the target calculation itself. However, because historical yield-class assignments come directly from Forestry Commission records, the original field assessments that generated them could not be independently verified.


\begin{table}[!htbp]
\centering
\scriptsize
\setlength{\tabcolsep}{3pt}
\renewcommand{\arraystretch}{0.95}
\begin{tabular}{
    >{\raggedright\arraybackslash}p{2.6cm}
    >{\raggedright\arraybackslash}p{7.0cm}
    >{\raggedright\arraybackslash}p{3.2cm}
}
\hline
Field & Meaning or derivation & Role in this study \\
\hline
\texttt{elev\_\allowbreak percentile\_\allowbreak 95th} & 95th percentile of ALS return heights above ground, used without correction & Response variable \\
\texttt{Top\_Height95} & $=\texttt{elev\_percentile\_95th}\times1.1$; Forest Research's own corrected height estimate & Retained for reporting but not modelled \\
\texttt{plyr}, \texttt{Age} & Planting year; \texttt{Age} = survey year $-$ \texttt{plyr} & Input feature; used for data cleaning \\
\texttt{yldc} & An existing Forest Research yield-class field, not calculated from this project's LiDAR-derived height response & Excluded from prediction inputs; used to construct Q2's own attribution target \\
\texttt{whcl} & Ordinal Forest Research windthrow hazard class, 0--6 & Management descriptor \\
\texttt{CanopyCover} & Canopy-cover fraction from the same ALS survey and processing step as the height response & Input feature \\
\texttt{Thin}, \texttt{last\_thinn}, \texttt{next\_thin\_} & Thinning history; \texttt{Thin} is binary (1 if thinned, 0 if never thinned); \texttt{last\_thinn}/\texttt{next\_thin\_} record last and next thinning years & Input feature \\
\texttt{Vol95}, \texttt{area}, \texttt{g1}, \texttt{g3} & Forest Research volume estimate, calculated from height, plot area, and canopy cover & Used for disturbance screening \\
\texttt{GYCspec95} & General Yield Class calculated directly from a plot's own measured height (\texttt{Top\_Height95}) & Excluded (target leakage) \\
\texttt{p1}--\texttt{p5} & Forest Research growth curve parameters for yield class & Excluded from inputs; used with \texttt{yldc} to build Q2 target curve \\
\hline
\end{tabular}
\caption{\small{Target, stand, and management fields provided by Forest Research. \texttt{GYCspec95} and \texttt{p1}--\texttt{p5} are excluded from model inputs to prevent target data leakage. In Q2, \texttt{p4} and \texttt{p5} are held fixed only to define the baseline yield-class asymptote target.}}
\label{tab:data-fields}
\end{table}

%------------------------------------------------
\section{Environmental data sources}
\label{sec:data-environmental}

Each plot was matched with external terrain, wind, climate, and soil variables (Table~\ref{tab:data-env-sources}). Appendix~\ref{tab:data-env-variables_detail_app} defines each variable.

\begin{table}[!htbp]
\centering
\scriptsize
\setlength{\tabcolsep}{4pt}
\renewcommand{\arraystretch}{1.05}
\begin{tabular}{ >{\raggedright\arraybackslash}p{2.0cm} >{\raggedright\arraybackslash}p{4.2cm} >{\raggedright\arraybackslash}p{2.3cm} >{\raggedright\arraybackslash}p{4.0cm}}
\hline
Data family & Source and resolution & Time basis & Variables used \\
\hline
Terrain & OS Terrain 50, 50\,m \cite{ordnancesurveyTerrain50_C3_LID} & Static & Elevation, slope, aspect, northness and eastness \\

Soil & CEH Natural Capital, 50\,m \cite{feeney2026CEHNaturalCapital_C3_LID}; SoilGrids v2.0, 250\,m \cite{poggio2021SoilGrids_C3_LID} & Static & Pedotope, drainage, texture and soil pH \\

Wind exposure & Global Wind Atlas, 250\,m \cite{davis2023GlobalWindAtlas_C3_LID} & 2006--2015 climate & Long-term wind speed and exposure \\

Observed wind & MIDAS Open weather stations \cite{metoffice2019MIDASOpen_C3_LID} & Survey intervals & Mean and extreme wind conditions \\

Long-term climate & CHELSA v2.1, 1\,km \cite{karger2021CHELSA_C3_LID} & 1981--2010 climate & Growing degree days and precipitation \\

Matched climate & HadUK-Grid, ~1\,km \cite{hollis2019HadUKGrid_C3_LID} & Cohort years & Temperature and frost \\

Spatial context & OS Open Roads and Open Rivers, vector data
\cite{ordnancesurveyOpenRoads_C3_LID,ordnancesurveyOpenRivers_C3_LID} & Static & Distance to roads, rivers and forest edges \\
\hline
\end{tabular}

\caption{\small{External data sources used in the analysis. Values were matched to plot centres. Full source, resolution, coverage and extraction details are given in Appendix~\ref{app:data-sources}.}}
\label{tab:data-env-sources}
\end{table}

Environmental variables vary in their time resolution. Terrain, soil, climate normals (CHELSA, Global Wind Atlas), and TOPEX were kept static, assigning each plot a single value across all survey years. TOPEX measures topographic shelter, where higher scores mean greater shelter and lower scores mean higher exposure. By contrast, MIDAS wind metrics were calculated specifically for each interval between survey dates using weather-station data. HadUK-Grid temperature and frost values were averaged over the survey years within each cohort, while road, river, and edge distances were treated as static baselines. Limited missing values in environmental features were handled by removing incomplete rows for specific model sets or using missing data indicators for thinning history (Chapter~\ref{ch:methodology}).

Several limitations should be noted: Requiring plots to appear across every survey year ensured clean tracking over time, but reduced landscape coverage, particularly in the 6-survey cohort where high-elevation areas are under-represented. Predictor layers also vary in detail: terrain changes across individual plots, whereas regional climate maps share the same values across wide areas and do not update per survey date. Additionally, canopy cover shares the same LiDAR flight acquisitions as top height, which could introduce shared survey conditions.
%=======================================================================================
%=======================================================================================

%=======================================================================================
%=======================================================================================
\chapter{Methodology and Experimental Framework} \label{ch:methodology}
\section{Study design}
\begin{table}[H]
\centering
\footnotesize
\renewcommand{\arraystretch}{1.1}
\setlength{\tabcolsep}{5pt}
\begin{tabularx}{\linewidth}{@{} l >{\hsize=0.85\hsize\raggedright\arraybackslash}X >{\hsize=0.85\hsize\raggedright\arraybackslash}X >{\hsize=1.3\hsize\raggedright\arraybackslash}X @{}}
\toprule
\textbf{Q} & \textbf{Unit \& Target} & \textbf{Models} & \textbf{Main Evaluation} \\
\midrule
Q1 & Plot; mean growth curve residual ($\bar r_i$) & Elastic Net, XGBoost, LMM & Spatial block CV, LMM variance decomposition, residual Moran's $I$ \\
\addlinespace[3pt]
Q2 & Plot; asymptote deviation from yield class ($\Delta y_{\max,i}$) & Elastic Net, XGBoost, GNNWR & 5-fold spatial block CV, SHAP attribution, Moran's $I$, target audit \\
\addlinespace[3pt]
Q3 & Plot--survey; top height ($H_{it}$) & CR, Linear, RF, XGBoost, DNN, PINN, PINN-$k$ & 5-fold spatial block CV, pooled metrics, hyperparameter/$\lambda$/feature checks, permutation importance, age/height errors \\
\bottomrule
\end{tabularx}
\caption{\footnotesize Units, targets, models, and evaluation methods across research questions.}
\label{tab:rq-framework}
\end{table}

\paragraph*{Data partitions and leakage control}
The 4-survey cohort was our main dataset. The smaller 6-survey cohort gave us a longer-term check. For Q1 and Q2, each plot became one summary value. For Q3, we used every plot-survey record. To stop data leakage, we split the data by forest compartment, not by plot \cite{roberts2017Cross_C5_EVAL}. No compartment appeared in both training and test sets. We also removed training plots within 60,m of validation or test plots, to stop spatial spillover. We grouped the compartments into five equal folds. In each round, three folds trained the model, one tuned its settings, and one tested final accuracy. Across the five rounds, each compartment was tested once. All test predictions were then pooled to get the overall scores \cite{ploton2020Spatial_C5_EVAL}.
  
\begin{figure}[!htbp]
\centering
\includegraphics[width=1.1\textwidth]{figures/fig_methodology/fig_method_map_splits.png}
\caption{Data-splitting strategies for model evaluation: 
(A)~Random plot-level split: shown for comparison only, as spatial proximity between adjacent training and test plots caused severe data leakage and inflated accuracy. 
(B)~Compartment-level spatial-block split with a 60\,m buffer: the primary strategy used throughout this dissertation to ensure complete spatial independence between training, validation, and test sets.}
\label{fig:splits-combined}
\end{figure}

Continuous variables were normalised using training-set statistics and applied to validation and test folds. Stand age was scaled separately so PINN outputs could be directly converted back to metres per year. Soil classes were one-hot encoded, while administrative IDs (compartment and block) were excluded from all feature sets. Additionally, two spatial neighbour-height variables were dropped because they were computed from adjacent plot heights prior to splitting and would have leaked test-set information.

%------------------------------------------------
\section{Feature-set construction}
Four feature sets were constructed to test whether adding environmental data improves model performance. Because feature relevance depends on the specific prediction task, sets were ranked and selected separately for Q1, Q2, and Q3. Variables with extensive missing values (e.g., $>20\%$) and exact duplicates were removed first. For highly correlated pairs (Spearman $|\rho| \geq 0.95$), we retained only the directly measured variable. The remaining candidates were ranked using their average rank across three metrics evaluated on a single spatial split: Spearman $|\rho|$, XGBoost permutation importance, and XGBoost drop-column ablation. To prevent multicollinearity, we iteratively removed variables with a variance inflation factor (VIF) greater than 5 until all remaining features met the threshold \cite{obrien2007Caution_C5_EVAL}. Baseline stand variables were retained unconditionally but still included in VIF calculations. This screening removed 430 four-survey plots, primarily due to gaps in SoilGrids coverage; the six-survey cohort had complete environmental coverage. Full missingness summaries and exact feature lists are provided in Appendix~\ref{app:feature-sets}.

\begin{table}[H]
\centering
\footnotesize
\renewcommand{\arraystretch}{1.2}
\begin{tabularx}{\linewidth}{@{}p{1.4cm}X@{}}
\toprule
Set & Construction \\
\midrule
Set~1 & Stand and management baseline only. \\
Set~2 & Set~1 plus the top 10 ranked environmental candidates (candidates failing VIF screening were replaced by the next-ranked variable). \\
Set~3 & Set~1 plus the top 50\% of terrain and wind features, followed by VIF screening. \\
Set~4 & Set~3 candidates plus the top 50\% of climate, soil, and spatial-context features, followed by VIF screening across the pooled set. \\
\bottomrule
\end{tabularx}
\caption{Construction of environmental feature sets: Feature ranking was performed independently for Q1, Q2, and Q3; full variable listings are given in Appendix~\ref{app:feature-sets}.}
\label{tab:feature-sets-construction}
\end{table}

Because VIF screening was rerun whenever the candidate pool expanded, the final sets are not strictly nested. For example, elevation was selected in Sets~2 and~3 but dropped from Set~4 due to collinearity with broader terrain and climate predictors. This reflects statistical redundancy within the larger pool rather than an absence of ecological influence.
%------
\section{Q1: Global departure attribution}
Q1 evaluates which forest stands consistently outperform or lag behind the regional (CR) growth curve once age is accounted for, and identifies the environmental drivers behind these differences. Repeated survey heights are condensed into a single mean residual per plot, and Elastic Net, XGBoost, and a linear mixed-effects model are trained to predict this deviation. We use the 4-survey cohort because the 6-survey cohort contains only 47 compartments, which is insufficient for robust spatial-block cross-validation.

\paragraph*{What Q1's target means in plain terms}
Let $H_{it}$ and $a_{it}$ denote the observed top height and stand age for plot $i$ at survey $t$, where $T_i$ is the total number of surveys for plot $i$. Using the regional baseline height $H(a_{it})$ from Equation~\eqref{eq:cr}, the survey-level residual $r_{it}$ and plot-level mean departure $\bar{r}_i$ are defined as:
\begin{align}
r_{it} &= H_{it} - H(a_{it}), \label{eq:cr-gap} \\
\bar{r}_i &= \frac{1}{T_i} \sum_{t=1}^{T_i} r_{it} \label{eq:mean-cr-residual}
\end{align}

A positive $\bar{r}_i$ indicates that a plot tracks above the regional baseline on average, while a negative value indicates under-performance. Because $\bar{r}_i$ reflects average vertical elevation relative to the curve rather than instantaneous growth rate, a plot's residual sign may vary across individual survey years. Stand age is excluded from the feature set since the baseline curve already removes the age trend.

\paragraph*{Models} Both Elastic Net and XGBoost are evaluated using 5-fold spatial-block cross-validation. Elastic Net tunes its regularisation penalty and $\ell_1$-ratio via inner 5-fold cross-validation (\texttt{ElasticNetCV}). XGBoost uses fixed hyperparameters (max depth 4, learning rate 0.04, up to 500 trees with early stopping), adopted after initial validation checks rather than exhaustive re-tuning on this specific target (the direct top-height XGBoost in Q3 uses its own targeted grid search). Model performance is reported using out-of-fold pooled metrics ($R^2$, RMSE, and MAE).

To account for plots grouped within compartments, a linear mixed-effects model is fit on each fold's training partition:
\begin{equation}
\bar{r}_i = \beta_0 + \boldsymbol{\beta}^{\top}\mathbf{x}_i + u_{c[i]} + \varepsilon_i, \qquad u_c \sim \mathcal{N}(0, \sigma_u^2).
\label{eq:mixed-effects}
\end{equation}
Here, $\boldsymbol{\beta}^{\top}\mathbf{x}_i$ is the fixed-effects term describing broad relationships across all plots (such as slope or TOPEX), and $u_{c[i]}$ is a random intercept for plot $i$'s compartment: each training compartment gets its own value, drawn from a shared distribution with variance $\sigma_u^2$, capturing grouping the model cannot otherwise explain (without saying why it exists). This model is used only for its fitted coefficients and the variance-explained comparison below, not to predict height for unseen compartments.

Parameters ($\boldsymbol{\beta}$, $\sigma_u^2$) are fitted with restricted maximum likelihood (REML), which avoids underestimating variance. Comparing a null model (compartment intercept only) against the full model (with fixed effects added) instead uses ordinary maximum likelihood, since REML fits cannot be compared across different fixed-effect structures; variance explained is the share of the null model's compartment-level variance that shrinks once fixed effects are added. The model is refit per fold, so coefficients are reported as mean $\pm$ SD across folds, showing fold-to-fold stability, unlike the pooled Elastic Net/XGBoost accuracy metrics above.

\paragraph*{Sensitivity checks} Canopy cover and top height come from the same LiDAR flights, so a strong canopy-cover effect could reflect measurement overlap rather than real stand structure. Three steps test this: refit Set~4 without canopy cover and check the change in $R^2$; check the direct correlation between canopy cover and each target; and recompute SHAP values without canopy cover to see if other predictors take over its role.

Terrain coefficients for slope and TOPEX also sometimes differ between Elastic Net and the mixed-effects model, which could mean local compartment effects that global models miss. A further check measures how much slope and TOPEX variance sits between compartments versus within them (the intraclass correlation coefficient, ICC).

Elastic Net coefficient signs, mixed-model fixed-effect signs, and XGBoost mean $|$SHAP$|$ rankings are then compared. This shows which variables matter, and in which direction when sign is available. The raw sizes are not directly comparable because the models differ, and mean $|$SHAP$|$ has no sign. Agreement across all three makes an association less likely to be a quirk of one model, but it does not show causality. Residual Moran's $I$ checks for remaining spatial clustering (Section~\ref{sec:eval-robustness}).
%------------------------------------------------
\section{Q2: Spatial asymptote attribution}
Q2 builds on Q1 by asking not just which variables matter, but whether their effect changes by place. Global models assume one environment-growth relationship for the whole study area; Q2 tests whether letting this relationship vary by location explains more and reduces leftover spatial clustering \cite{du2020Geographically_C6_GWR}. Each plot gets a prediction based on its location, on the idea that nearby plots behave more alike than distant ones, a neural network learns these location-based weights directly from coordinates, with no fixed neighbourhood size chosen in advance. GNNWR still assumes each local relationship is a straight line, so a changing local coefficient might just mean the straight-line fit is poor there, not a real difference in growth drivers by place. For that reason, both accuracy and the local coefficients are considered together.


\paragraph*{What Q2's target means in plain terms} Q2 asks whether a plot grows towards a higher or lower final height than its official yield class predicts. $p_1$--$p_5$ are Forest Research's published numbers that convert a yield-class rating into a full growth curve for a species and region. This project did not derive them, just reused them. $p_1$--$p_3$ set the curve's asymptote (ceiling); $p_4$--$p_5$ set its shape (rate and curvature), two parts of one sigmoid. $p_4$ and $p_5$ are fixed at their published values for the matching Forest Research curve; only the height ceiling is allowed to differ from plot to plot. Earlier work has represented environmental variation through the Chapman--Richards asymptote in the same way \cite{scolforo2013Dominant_C2_CR}. This supports it as a simple and workable choice. For plot $i$ at survey $t$:
\begin{align}
g_{it} &= \left(1-e^{-p_{4,i}a_{it}}\right)^{p_{5,i}}, \label{eq:git-curve} \\
\hat y_{\max,i} &= \frac{\sum_t H_{it}g_{it}}{\sum_t g_{it}^2}, \label{eq:plot-ymax-fit} \\
y_{\max,i}^{\mathrm{yldc}} &= \mathrm{yldc}_i p_{2,i}+p_{1,i}+2p_{3,i}, \label{eq:yldc-asymptote} \\
\Delta y_{\max,i} &= \hat y_{\max,i}-y_{\max,i}^{\mathrm{yldc}}. \label{eq:local-ymax-diff}
\end{align}
With $p_4$ and $p_5$ fixed, Equation~\eqref{eq:plot-ymax-fit} gives a least-squares estimate for $\hat y_{\max,i}$. This avoids trying to estimate three correlated curve parameters from only 4--6 observations per plot. $\hat y_{\max,i}$ is the ceiling fitted from plot $i$'s own observed trajectory, $y_{\max,i}^{\mathrm{yldc}}$ (Equation~\eqref{eq:yldc-asymptote}) is the ceiling implied by its assigned yield class, using the same fixed rate and shape. Their difference, $\Delta y_{\max,i}$ (Equation~\eqref{eq:local-ymax-diff}), is positive when the fitted ceiling is higher than expected and negative when it is lower. This target is estimated, not directly observed.

Following the Q1 setup, Elastic Net and XGBoost are trained on $\Delta y_{\max}$ using Sets~2--4, with the same tuning approach as Q1. Elastic Net tests a global linear relationship. XGBoost tests a global non-linear relationship. GNNWR tests a spatially varying but locally linear relationship. Spatial variation is modelled using GNNWR via its Python library \cite{du2020Geographically_C6_GWR, du2022PyGNNWR_C6_GWR}:
\begin{equation}
\widehat{\Delta y}_{\max,i} = \beta_0(s_i) + \sum_j \beta_j(s_i) x_{ij}.
\label{eq:gnnwr}
\end{equation}
For each location $s_i$, GNNWR uses a neural network to learn spatial weights from distances to training plots. These weights determine the local coefficients $\beta_j(s_i)$ around a global linear baseline. Each training fold uses only that fold's own training partition. No hyperparameter search was performed for GNNWR. Every run uses the library's default architecture and training settings.


\subsection*{Computational adjustments and feature evaluation}
Headline GNNWR runs used full reference sets (earlier GPU-memory caps are retained only for sensitivity checks). To bypass memory limits, hat-matrix diagnostics were disabled; consequently, GNNWR AIC uses raw feature counts rather than effective degrees of freedom, and $F$-tests are omitted (predictive metrics, Elastic Net, and XGBoost are unaffected). Feature importance is evaluated via XGBoost gain, test-fold SHAP values \cite{lundberg2017Unified_C7_XAI}, and GNNWR local coefficient distributions and maps. Feature-group contributions to $R^2$ use permutation importance (shuffling columns without refitting), while residual clustering impacts are measured by removing groups from Set~4, refitting, and comparing Moran's $I$ on the four-survey cohort. Finally, canopy cover dependency is isolated by refitting all three models using canopy cover exclusively and with it removed.

\paragraph*{Target and outlier audit}
Plots affected by harvesting or clear measurement errors are screened using two rules across consecutive surveys: clearfelling requires drops in top height $\ge 25\%$, canopy cover $\ge 70\%$, and volume $\ge 70\%$, while measurement error is flagged when top height drops $\ge 25\%$ alongside canopy cover drops $\le 10\%$ and stable volume; unclassified height drops are retained. Extreme fitted height ceilings ($\hat{y}_{\max,i}$) are flagged using Tukey's $1.5 \times \text{IQR}$ fences \cite{tukey1977Exploratory_C5_EVAL},
\begin{equation}
\left[Q_1 - 1.5 \times \text{IQR},\; Q_3 + 1.5 \times \text{IQR}\right],
\label{eq:tukey-fences}
\end{equation}
These bounds are calculated per cohort, Flagged plots remain in main models but are evaluated separately in a holdout check. To determine if extreme errors link to fitting artefacts, disturbance, or environment, we compare overlap among the top 1\% and 5\% residual plots across Elastic Net, XGBoost, and GNNWR, examining their spatial spread, edge distance, wind shelter, and local coefficients. Results are detailed later.

%------------------------------------------------
\section{Q3: Top-height prediction}
Q3 compares models for predicting top height: a baseline growth curve, standard tabular models, and neural networks, including physics-informed ones. It then tests whether environmental variables or Chapman--Richards(CR) guidance improve accuracy or the interpretability of the fitted parameters. Using five spatial folds on Set~3, the other feature sets test whether adding environmental data helps.

\paragraph*{CR baseline}
\begin{equation}
H(a) = y_{\max}\left(1 - e^{-ka}\right)^p,
\label{eq:cr}
\end{equation}
where $H(a)$ is height at age $a$, $y_{\max}$ is the height ceiling, $k$ is the growth rate, and $p$ controls curve shape. All three parameters are fit freely by non-linear least squares on each training fold and then frozen, keeping test folds fully independent.

\paragraph*{Conventional models}
We test three baselines: linear regression, random forest, and XGBoost. For Q3, XGBoost was tuned across 27 parameter settings over five folds using validation $R^2$ (Appendix~\ref{app:hyperparameters}).

\paragraph*{DNN and physics-informed models}
All neural networks: three-layer trunk, 128 units per layer, Leaky-ReLU, Adam optimiser (learning rate $10^{-4}$, weight decay $10^{-5}$, batch size 256, early stopping), defaults, checked against alternatives (Sensitivity checks and diagnostics, below). Standard DNN: minimises mean squared error ($\mathcal{L}_{\mathrm{data}} = \frac{1}{N}\sum (\hat H_i-H_i)^2$). Conditioned DNN: environmental variables feed the input directly. PINN, PINN-$k$: environmental variables feed small sub-networks instead, which adjust each plot's own growth-curve parameters. Plain PINN adjusts only the height ceiling ($y_{\max}$); growth rate ($k$) and shape ($p$) stay one shared value for every plot. PINN-$k$ adjusts both $y_{\max}$ and $k$; $p$ stays shared for both variants. These plot-level parameters adjust the predictions and loss terms together, balancing fit to the data with the baseline growth curve.

Equations~\eqref{eq:cr-derivative}--\eqref{eq:total-loss}: the physics-informed loss.
\begin{align}
H'(a) &= y_{\max}pk e^{-ka}\left(1-e^{-ka}\right)^{p-1}, \label{eq:cr-derivative} \\
\mathcal{L}_{\mathrm{phys}} &= \frac{1}{N}\sum_{i=1}^{N}\left(\frac{\partial\hat H_i}{\partial a_i}-H'(a_i)\right)^2, \label{eq:physics-loss} \\
\mathcal{L}_{\mathrm{traj}} &= \frac{1}{M}\sum_{j=1}^{M}\left(\frac{\hat H_{j,\mathrm{later}}-\hat H_{j,\mathrm{earlier}}}{\Delta a_j}-H'(\bar a_j)\right)^2, \label{eq:trajectory-loss} \\
\mathcal{L} &= \mathcal{L}_{\mathrm{data}} + \lambda_{\mathrm{phys}}\mathcal{L}_{\mathrm{phys}} + \lambda_{\mathrm{traj}}\mathcal{L}_{\mathrm{traj}} + 10^{-5}\sum_{\theta}|\theta|. \label{eq:total-loss}
\end{align}
Equation~\eqref{eq:cr-derivative}: theoretical growth rate, shown with the shared global parameters; conditioned models substitute each plot's own adjusted parameters instead, $y_{\max}^{(i)}$ for both PINN variants, plus $k^{(i)}$ for PINN-$k$. Equation~\eqref{eq:physics-loss}: gap between this rate and the network's own derivative ($\partial\hat H_i/\partial a_i$). Equation~\eqref{eq:trajectory-loss}: gap between consecutive-survey growth and this rate, at midpoint age $\bar a_j$. Equation~\eqref{eq:total-loss}: data loss plus weighted physics and trajectory terms, plus an L1 penalty. Defaults: $\lambda_{\mathrm{phys}}=\lambda_{\mathrm{traj}}=1$; 0 turns the physics terms off.

Sub-networks adjust plot parameters:
\begin{equation}
y_{\max}^{(i)} = y_{\max}+f_y(\mathbf{x}_i), \quad k^{(i)} = k\exp\!\left(f_k(\mathbf{x}_i)\right).
\label{eq:param-adjust}
\end{equation}
Plain PINN: $y_{\max}^{(i)}$ line only, $k^{(i)}=k$ fixed. PINN-$k$: both lines. Each sub-network in use ($f_y$ alone for PINN; $f_y$, $f_k$ for PINN-$k$): one hidden layer, 16 units, turns terrain features into one number. Trunk and sub-networks train together: one Adam optimiser, one learning rate, one weight decay, no separate tuning step per part. Sub-network outputs feed the prediction directly (Equation~\eqref{eq:pinn-prediction}), so ordinary training already shapes them; $\lambda_{\mathrm{phys}}$ adds a further pull towards matching the theoretical growth rate.

The final prediction: the CR curve at the plot's adjusted parameters, plus a residual from the trunk:
\begin{equation}
\hat H_i = y_{\max}^{(i)}\left(1-e^{-k^{(i)}a_i}\right)^{p} + \text{trunk residual}.
\label{eq:pinn-prediction}
\end{equation}
$y_{\max}^{(i)}$ shapes the prediction directly for both variants; $k^{(i)}$ does too, for PINN-$k$.

\subsection*{Sensitivity checks and diagnostics}
Beyond 5-fold spatial evaluation, six targeted diagnostics test the robustness of the PINN and PINN-$k$ models (single run per fold/condition):
\begin{itemize}
\setlength\itemsep{1pt}
\item Training and physics settings: Evaluated alternative learning rates, batch sizes (512 vs.\ 256), and eight physics weights ($\lambda_{\mathrm{phys}} = \lambda_{\mathrm{traj}} \in [0, 3.0]$) across all models.
\item Feature-set sensitivity: Tested models without environmental inputs and across Sets~2, 3, and 4 to assess input necessity and feature selection sensitivity.
\item Terrain permutation importance: Shuffled individual terrain/wind inputs (10 repeats) to quantify their impact on personalised $y_{\max}$ and $k$ predictions.
\item Age/height stratification: Compared MAE across age and height bands to evaluate error distributions at juvenile and mature extremes.
\item Flagged compartment tracking: Monitored a large-departure compartment across all five folds (in-sample vs.\ out-of-sample) to separate true site signals from fold-specific overfitting.
\item Yield-class alignment: Correlated predicted parameter deviations against withheld Forest Research yield classes across the full test set to verify biological realism.
\end{itemize}


\section{Evaluation and robustness} \label{sec:eval-robustness}
Model accuracy is evaluated using RMSE, MAE, $R^2$, and Bias. Here $e_i = y_i - \hat y_i$, where $y_i$ is the relevant target: $\bar r_i$ for Q1, $\Delta y_{\max,i}$ for Q2, and top height for Q3.
\begin{equation}
\mathrm{RMSE}=\sqrt{\tfrac{1}{N}\textstyle\sum e_i^2}, \quad
\mathrm{MAE}=\tfrac{1}{N}\textstyle\sum |e_i|, \quad
R^2=1-\tfrac{\sum e_i^2}{\sum (y_i-\bar y)^2}, \quad
\mathrm{Bias}=\tfrac{1}{N}\textstyle\sum e_i.
\label{eq:eval-summary}
\end{equation}
Positive bias indicates under-prediction. We report fold-level means ($\pm\text{SD}$) or pooled metrics across all out-of-fold predictions. For the 4-survey cohort, $R^2$ $95\%$ confidence intervals are calculated using 1,000 compartment bootstrap resamples.

Residual spatial autocorrelation is measured using global Moran's $I$:
\begin{equation}
I=\frac{N}{W}\frac{\sum_i\sum_j w_{ij}(e_i-\bar e)(e_j-\bar e)}{\sum_i(e_i-\bar e)^2}.
\label{eq:morans-i}
\end{equation}
Spatial weights ($w_{ij}$) use binary, row-standardised distance bands matched to each model's own semivariogram range (up to 5,000\,m; 999 permutations). When needed, rows are subsampled to 5,000 to keep this fast to compute. Moran's $I$ is left out if semivariance (the difference between pairs of points at increasing distances) is flat. It is marked provisional if semivariance keeps rising past 5,000\,m. Distances are reported alongside every value. Because each model's band is fitted to its own residual surface, Moran's $I$ values are comparable in direction across models, but not as exact like-for-like magnitudes.

Seed and split sensitivity checks were applied selectively based on computational limits. For Elastic Net and XGBoost, models were run across eight spatial fold splits for both cohorts to test sensitivity to boundary definitions. GNNWR was tested across three splits on the 6-survey cohort (omitting random training seeds), while DNN, PINN, and PINN-$k$ used a single seed per fold, relying instead on architectural and physics-weight tests for stability evaluation. For PINN-$k$, parameter disentanglement between final height ($y_{\max}$) and growth rate ($k$) was assessed via correlation across 2,000 compartment bootstraps, reported alongside $R^2$ for the main setup and physics-weight tests.

%=======================================================================================
%=======================================================================================
\chapter{Results and Discussion} \label{ch:results}

\section{Q1: Global departure attribution}
Using the 4-survey cohort, canopy cover and thinning history dominate the attribution in every feature set and every method tested (Table~\ref{tab:results-q1}).
\begin{table}[H]
\centering
\scriptsize
\renewcommand{\arraystretch}{1.2}
\begin{tabular*}{\textwidth}{@{\extracolsep{\fill}}lrrrr@{}}
\toprule
Set & n plots & EN $R^2$ [95\% CI] & XGB $R^2$ [95\% CI] & LMM spatial var. explained \\
\midrule
Set1 (baseline only) & 71{,}766 & 0.220 [0.180, 0.252] & 0.222 [0.183, 0.252] & 0.016$\pm$0.078 \\
Set2 & 71{,}766 & 0.359 [0.310, 0.396] & 0.416 [0.355, 0.461] & 0.204$\pm$0.066 \\
Set3 & 71{,}727 & 0.325 [0.272, 0.365] & 0.375 [0.321, 0.414] & 0.049$\pm$0.048 \\
Set4 & 71{,}330 & 0.358 [0.305, 0.398] & 0.395 [0.338, 0.438] & 0.175$\pm$0.064 \\
\bottomrule
\end{tabular*}
\caption{\footnotesize{Attribution of the persistent mean curve residual (Equation~\ref{eq:mean-cr-residual}),
4-survey cohort, spatial-block split. EN = Elastic Net; XGB = XGBoost; LMM = linear mixed-effects model.
Set1 = canopy cover + thinning only (baseline); Set2 = top-10 ranked environmental predictors; Set3 =
terrain/wind only; Set4 = full environmental set (terrain + climate + soil). $R^2$ 95\% CIs from 1,000
compartment bootstrap resamples. LMM spatial variance explained = share of variance attributable to the
compartment random intercept on top of that set's fixed effects (mean $\pm$ SD across folds, not a
bootstrap CI).}}
\label{tab:results-q1}
\end{table}

These $R^2$ scores are modest: even the best model (Set~2, XGBoost, $R^2=0.416$) leaves over half the growth variance unexplained. The goal is to uncover key environmental drivers rather than make exact plot predictions. Note that $R^2$ is not directly comparable across Q1, Q2, and Q3 because each uses a different target with its own variance. Figure~\ref{fig:q1-fig4-trajectories} illustrates these persistent curve departures in the observed data across six example plots spanning the residual range.

\begin{figure}[!htbp]
  \includegraphics[width=1.1\textwidth]{figures/fig_results/q1/q1_final_fig4_example_trajectories.png}
  \centering
 \caption{\footnotesize{Example plots by residual tier. Observed top height vs.\ age
  (solid) against the pooled Chapman--Richards curve (dashed), for five 4-survey plots: two near
  the 90th percentile of $\bar r_i$ (above curve), one near zero (typical), two near the 10th
  percentile (below curve).}}
\label{fig:q1-fig4-trajectories}
\end{figure}

\subsection{Discussion}
%-----
\paragraph*{Canopy cover and thinning history dominate attribution across all models and feature sets.}
Canopy cover has the largest positive Elastic Net coefficient (Set~4: $+1.943 \pm 0.034$) and the top XGBoost importance ($|\text{SHAP}| = 1.375\text{--}1.442$). Since both canopy cover and the target come from the same LiDAR survey, we ran three checks to see if it is just a proxy for the target itself (Figure~\ref{fig:q1-fig1-canopycover-ablation}). Only moderate correlation ($r = 0.35\text{--}0.47$) rules out a simple duplicate. Removing canopy cover drops $R^2$ by roughly a third (Elastic Net: $0.350 \to 0.231$; XGBoost: $0.388 \to 0.249$; Figure~\ref{fig:q1-fig1-canopycover-ablation}, left), showing it adds real, unique information. No other variable steps in to fill the gap either; $\mathit{dist}_{\text{to\_road}}$'s importance barely moves ($|\text{SHAP}| = 1.003 \to 0.869$; Figure~\ref{fig:q1-fig1-canopycover-ablation}, right). So canopy cover and thinning history are treated as a structural baseline, and the real attribution question is what else adds on top of them.

\begin{figure}[!htbp]
\centering
 \includegraphics[width=\textwidth]{figures/fig_results/q1/q1_final_fig1_canopycover_ablation.png}
 \caption{\footnotesize{CanopyCover carries real, non-redundant signal. Left: pooled $R^2$ (Set4, 4-survey) with vs.\ without CanopyCover for Elastic Net and XGBoost, both lose roughly a third of their explanatory power. Right: mean $|$SHAP$|$ for CanopyCover, $\mathit{dist\_to\_road}$, precipitation and slope, with vs.\ without CanopyCover, $\mathit{dist\_to\_road}$'s importance does not rise to compensate, evidence against simple correlation-driven credit transfer.}}
\label{fig:q1-fig1-canopycover-ablation}
\end{figure}

%-----
\paragraph*{Most Set~4 predictors agree across models, but six exhibit sharp divergences} anchored by $\mathit{topex}$ (Figure~\ref{fig:q1-fig2-attribution-lenses}; Appendix~\ref{app:nonlinearity-check}). While $\mathit{slope\_degrees}$ and rainfall remain consistently positive across LMM and Elastic Net (slope: $+0.86\text{-}+1.23$), $\mathit{topex}$ splits sharply (LMM $+1.32$ vs.\ Elastic Net $-0.18$ to $+0.06$). This is not driven by multicollinearity (unstable predictors have lower VIFs, 2.1--3.4, than stable rainfall's 3.91) or climate-data timing (the static rainfall normal is stable, while the survey-matched $\mathit{tas\_mean}$ is unstable). All unstable predictors show high between-compartment clustering ($\text{ICC} = 0.74$--$0.89$), yet stable rainfall's ICC is even higher ($0.97$), so clustering alone cannot explain the divergence, the real mechanism follows below.
%-----

\paragraph*{Within-compartment variation separates genuine local signals from spatial artefacts.}
Because both models assume linear relationships, their divergence stems from spatial handling: LMM uses random compartment intercepts, whereas Elastic Net pools landscape-wide data. To find why coefficients disagree, we calculated within-compartment correlations by mean-centering predictors and targets per compartment, removing regional gradients to test if variables drive growth differences between neighbours.

This check sorts the six variables with sign disagreements (Figure~\ref{fig:q1-fig2-attribution-lenses}, left; Appendix~\ref{app:nonlinearity-check}). Only one, $\mathit{topex}$ (green ring), retains a true local link: its correlation strengthens after removing compartment differences ($0.08 \to 0.21$), letting LMM recover a physical effect that pooled models obscure. Conversely, the other five (red rings) see correlations shrink towards zero—wind speed ($\mathit{gwa\_weibull\_a\_50m}$, $-0.33 \to -0.14$), temperature ($\mathit{tas\_mean}$, $0.28 \to 0.09$), soil pH ($\mathit{soilgrids\_ph}$, $0.22 \to 0.04$), wetness ($\mathit{ceh\_twi}$, $-0.11 \to -0.01$), and local relief ($\mathit{local\_relief\_500m}$, $-0.12 \to -0.01$)—showing raw correlations reflect compartment placement rather than local effects. Rainfall ($\mathit{chelsa\_bio12\_precip\_mm}$, unringed) maintains agreement between models. Predictor instability thus stems from spatial confounding, not unmodelled nonlinearity.

\begin{figure}[H]
\centering
\includegraphics[width=\textwidth]{figures/fig_results/q1/q1_final_fig2_attribution_three_lenses.png}
\caption{\footnotesize{Comparing how different models rank and score environmental factors (Set~4, 4-survey cohort). Left: Model scores (average and range across 5 folds) for Elastic Net (purple) and Mixed-Effects (grey). Set~3 is shown with hollow markers, Set~4 with solid ones. For the six variables where the two models disagreed, green rings mark variables with a true local effect like terrain shelter, while red rings mark variables whose apparent effect disappears when looking within individual compartments. Rainfall is left unringed as a stable baseline. Right: XGBoost importance scores ($\text{SHAP}$) for the same variables.}}
\label{fig:q1-fig2-attribution-lenses}
\end{figure} 
%-----
\paragraph*{Spatial clustering falls unevenly, and never fully disappears}
Every feature set still leaves real spatial clustering in its prediction errors (distance-band Moran's $I$ at 1385--1739\,m, $p = 0.001$; Figure~\ref{fig:q1-fig3-spatial-clustering}B). Going from the stand-structure baseline (Set~1) to the full environmental model (Set~4), Elastic Net cuts this clustering by 29\% ($0.145 \to 0.103$) and XGBoost by 63\% ($0.146 \to 0.054$).

But XGBoost's drop is not smooth: clustering falls sharply at Set~2 ($0.146 \to 0.032$, a 78\% drop), bounces back up at Set~3 ($0.077$), then falls again at Set~4 ($0.054$). This bounce comes from road distance ($\mathit{dist\_to\_road}$) and macroclimate variables ($\mathit{chelsa\_bio12\_precip\_mm}$, $\mathit{tas\_mean}$), which are missing from Set~3. Dropping these from Set~4 pushes clustering back up to $0.084$ but costs almost no accuracy ($R^2$ only falls from $0.395$ to $0.368$). The clustering left in both models points to environmental drivers we have not measured, or spatial patterns that simply vary by place \cite{gaspard2019Residual_C5_EVAL}.

\begin{figure}[H]
\centering
  \includegraphics[width=\textwidth]{figures/fig_results/q1/q1_final_fig3_spatial_clustering.png}
\caption{\footnotesize{Spatial clustering of model residuals (Set~4, 4-survey cohort). 
  Panel A (top): Set~4 Elastic Net residuals averaged by compartment, showing leftover spatial errors examined in Q2. Panel B (bottom): Residual Moran's $I$ across feature sets. XGBoost reduces clustering more than Elastic Net overall, but drops sharply at Set~2 before bouncing back at Set~3.}}
\label{fig:q1-fig3-spatial-clustering}
\end{figure}

Set~3 (terrain-only) is the weakest option. Refitting Set~4 category by category shows soil adds only a small gain ($+0.022$ $R^2$), while macroclimate actually hurts accuracy ($-0.056$ $R^2$; Appendix~\ref{app:category-ablation}). Terrain and edge-effect features drive most of Set~4's performance instead. This goes against the common assumption that forest yield modelling needs broad climate layers \cite{bontemps2014Predictive_C2_BG}, here, regional climate grids add noise, not signal.

Every environmental model clearly beats the stand-structure baseline (Set~1), but XGBoost's confidence intervals for Set~2 and Set~4 overlap heavily ($[0.355, 0.461]$ vs.\ $[0.338, 0.438]$). So we cannot say with confidence that the full environmental set (Set~4) is really better than the simpler Set~2.


%-----

\subsection*{Q1: Conclusion}
Stand structure and management history account for most plot-level growth variations. Among environmental factors, slope and rainfall remain stable across models, but six other variables show major disagreements between Elastic Net and the mixed-effects model. This happens because Elastic Net pools data across the whole landscape, turning regional artefacts into false signals; the mixed-effects model avoids this by looking within individual compartments, recovering true local effects for features like $\mathit{topex}$ and wind exposure while exposing the others as regional noise. Overall, soil properties add almost nothing to predictions ($+0.022$ $R^2$), and macroclimate grids actively hurt accuracy ($-0.056$ $R^2$). Finally, significant spatial errors persist in all models' residuals—meaning broad environmental data cannot fully explain regional patterns on its own. This leftover spatial error directly motivates the spatially varying coefficient models evaluated in Q2.

%=======================================================================================
\section{Q2: Spatial asymptote attribution}
Q1 showed that global models leave strong spatial clustering in their prediction errors. Q2 tests whether a spatially varying model (GNNWR) can capture what these global models miss.

In terms of mean accuracy, GNNWR leads across every feature set (Table~\ref{tab:results-q2-ablation}). On Set~4, it reaches $R^2 = 0.294$ ($[0.236, 0.347]$), ahead of Elastic Net ($0.240$, $[0.192, 0.286]$) and XGBoost ($0.250$, $[0.201, 0.295]$). However, broad overlap in confidence intervals means this lead is not clear-cut, and even its best score ($R^2 = 0.320$ on Set~2) leaves most growth variance unexplained. Looking into GNNWR's small advantage shows that because all three models rely almost entirely on canopy cover, GNNWR gains its edge simply by locally smoothing this single main feature, rather than finding subtle environmental trends missed by global baselines.

\begin{table}[h]
\centering
\footnotesize
\renewcommand{\arraystretch}{1.2}
\begin{tabular*}{\textwidth}{@{\extracolsep{\fill}}lrrr@{}}
\toprule
Method & Full Set4 & No CanopyCover & CanopyCover only \\
\midrule
Elastic Net & 0.240 (0.223$\pm$0.100) & 0.042 (0.024$\pm$0.077) & 0.207 (0.190$\pm$0.084) \\
XGBoost & 0.250 (0.231$\pm$0.119) & 0.061 (0.040$\pm$0.094) & 0.210 (0.193$\pm$0.081) \\
GNNWR & 0.294 (0.277$\pm$0.068) & 0.049 (0.026$\pm$0.067) & 0.249 (0.235$\pm$0.030) \\
\bottomrule
\end{tabular*}
\caption{\small{Impact of canopy cover on model performance (Set~4, 4-survey cohort): Pooled $R^2$ (5-fold mean $\pm$ SD in parentheses). All three models collapse near zero without canopy cover ($R^2 < 0.07$). Canopy cover alone accounts for $\sim$80--83\% of each model's predictive power across both global baselines (Elastic Net, XGBoost) and the spatially-varying model (GNNWR).}}
\label{tab:results-q2-ablation}
\end{table}

\subsection{Discussion}
\paragraph*{A noisy target explains the weak signal elsewhere}
Part of the reason other variables barely register is that Q2's target contains heavy noise. The fitting step fixes each plot's growth rate to its assigned yield class and estimates only the maximum height. When a plot grows faster or slower than assumed, that difference does not vanish; it spills straight into the estimated ceiling.
\begin{figure}[H]
\centering
\includegraphics[width=\textwidth]{figures/fig_results/q2/q2_example_plot_curve.png}
\caption{\small{Example plot-level asymptote gap ($\Delta y_{\max}$) and XGBoost attribution. Left: Observed trajectory for Plot 309770, its fitted CR curve, and the yield-class reference curve, with the shaded region showing the observed gap across survey ages 65--81. Right: Local SHAP feature contributions for the same plot, illustrating how individual predictors account for the estimated ceiling departure.}}
\label{fig:q2-example-plot}
\end{figure}

Figure~\ref{fig:q2-example-plot} shows this for Plot 309770, chosen because it is unusually well-explained, not typical: its fitted ceiling reaches 43.0\,m, well above its reference value of 32.8\,m (a $+10.2$\,m gap), and XGBoost predicts $+12.4$\,m, with canopy cover actually lowering the predicted ceiling here, unlike its usual positive effect. A typical plot, picked by median prediction error (Appendix Figure~\ref{fig:q2-example-plot-typical}), tells a different story: the model captures only about half the observed gap, with no single factor standing out, making this clean example the exception. Because $\Delta y_{\max}$ merges ceiling error and growth-rate mismatch into one number, only very strong signals like canopy cover stand out clearly, weaker, real environmental effects get lost, which is why most other site variables show little impact.

\paragraph*{Q1 and Q2 agree on which variables to trust}
Q1 asks whether Elastic Net and the mixed-effects model (LMM) agree on the direction of a variable's effect (positive or negative). Q2 asks whether GNNWR's local coefficient points the same direction in every forest compartment, or flips between compartments. Canopy cover and slope pass both checks cleanly: Elastic Net and LMM agree on their direction in Q1, and in Q2 their GNNWR coefficients point the same way in almost every compartment (canopy cover 100\%, slope 97\% of the 231 compartments). Rainfall agrees on direction in Q1 too, and is mostly one-directional in Q2 (87\% of compartments), though less uniformly so than canopy cover or slope. Soil pH fails both checks: Elastic Net and LMM disagree on its direction in Q1, and in Q2 its GNNWR coefficient is positive in barely half the compartments (51\%) and negative in the rest, meaning no real, dependable relationship at either scale (Appendix Table~\ref{tab:results-q2-reliability}).
 
Topographic exposure ($\mathit{topex}$) is the interesting middle case. In Q1, Elastic Net and LMM disagree on its direction because $\mathit{topex}$ has a real, fairly strong effect \emph{within} each compartment (correlation $0.214$) that looks much weaker once compartments are pooled together (correlation $0.080$), so a model that only sees the pooled data gets the direction wrong. Q2 shows the same thing from a different angle: GNNWR's local coefficient for $\mathit{topex}$ is positive in 72\% of compartments but negative in the other 28\% (Figure~\ref{fig:q2-stability-map}). Because Q1 and Q2 use entirely different statistical methods and still tell the same story, this instability looks like a real feature of the landscape, not an artefact of one particular model.

\begin{figure}[!htbp]
\centering
\includegraphics[width=\textwidth]{figures/fig_results/q2/q2_topex_vs_canopycover_stability.png}
\caption{\small{Spatial stability of GNNWR's local coefficients (Set~4, 4-survey cohort). Each map shows the compartment-average local coefficient on the same colour scale: how much $\Delta y_{\max}$ (m) changes per one-unit increase in the predictor. Left: canopy cover's coefficient stays positive in every compartment. Right: topographic exposure's ($\mathit{topex}$) coefficient is positive in most compartments but flips negative in over a quarter of them. This is why Elastic Net and LMM (Q1) disagree on $\mathit{topex}$'s direction: a global model forced to pick one sign for the whole landscape cannot capture this compartment-by-compartment flip.}}
\label{fig:q2-stability-map}
\end{figure}

A few minor limitations are worth noting. Some terrain features (like wind exposure and slope) show nonlinear patterns—rising before levelling off—which GNNWR's locally linear setup cannot fully capture. However, because canopy cover drives 80--90\% of explained variance, this unmodelled nonlinearity has little practical impact. A tiny fraction of stands (under 1\% of the 4-survey cohort; 266 of 56,526 plots) show extreme trade-offs between fitted growth rates and maximum heights (Appendix~\ref{app:flagged-plot-split}). Finally, the smaller 6-survey cohort produced near-zero accuracy across all models, reflecting its far smaller spatial sample (47 vs.\ 231 compartments).

\subsection*{Q2: Conclusion}
Q2's result is a useful negative finding: letting the model vary its coefficients across space (GNNWR) did not reveal any hidden environmental driver beyond canopy cover, which the simpler global models had already found on their own. The diagnostics explain why. The target itself mixes two different kinds of error into one number, how fast a tree grows, and how tall it eventually gets, which makes any clean spatial pattern harder to find. GNNWR can only fit straight-line relationships within each local area, so it cannot capture terrain effects that curve or level off. In short, these are limits of the method and the target, not evidence that spatial ecology does not matter here.

%=======================================================================================
\section{Q3: Top-height prediction}

% --- Table 1: six-model comparison, 4-survey only ---
\begin{table}[H]
\centering
\scriptsize
\renewcommand{\arraystretch}{0.9}
\setlength{\tabcolsep}{12pt}
\begin{tabular}{l ccc}
\toprule
\textbf{Model} & \textbf{RMSE (m)} $\downarrow$ & \textbf{MAE (m)} $\downarrow$ & $\mathbf{R^2}$ $\uparrow$ \\
\midrule
Linear          & $5.168 \pm 0.361$ & $4.056 \pm 0.343$ & $0.580 \pm 0.026$ \\
RF              & $4.794 \pm 0.187$ & $3.611 \pm 0.142$ & $0.638 \pm 0.011$ \\
XGBoost (tuned) & $\mathbf{4.548 \pm 0.137}$ & $\mathbf{3.431 \pm 0.123}$ & $\mathbf{0.674 \pm 0.009}$ \\
DNN             & $4.681 \pm 0.230$ & $3.533 \pm 0.201$ & $0.655 \pm 0.016$ \\
PINN$^*$        & $4.842 \pm 0.321$ & $3.689 \pm 0.289$ & $0.631 \pm 0.027$ \\
PINN-$k^*$      & $4.927 \pm 0.302$ & $3.745 \pm 0.271$ & $0.618 \pm 0.023$ \\
\bottomrule
\end{tabular}
\caption{\footnotesize{Top-height model performance on Set~3 (4-survey cohort; mean $\pm$ SD across five spatial folds; $N=232{,}292$ row-years across $58{,}073$ plots). Bold indicates best performance. $^*$PINN and PINN-$k$ link predicted parameters directly to height predictions and are evaluated without hyperparameter tuning. Results for the 6-survey cohort are in Appendix~\ref{app:results-q3-6survey}.}}
\label{tab:results-q3}
\end{table}

% --- Table 2: PINN/PINN-k across feature sets and hyperparameter checks ---
\begin{table}[H]
\centering
\scriptsize
\renewcommand{\arraystretch}{0.9}
\setlength{\tabcolsep}{18pt}
\begin{tabular}{lcc}
\toprule
\textbf{Condition} & \textbf{PINN} $\mathbf{R^2}$ & \textbf{PINN-}$\mathbf{k}$ $\mathbf{R^2}$ \\
\midrule
\multicolumn{3}{l}{\textit{Feature set}} \\
No environment               & $0.572 \pm 0.032$ & $0.575 \pm 0.033$ \\
Set~2 (small)                & $0.627 \pm 0.023$ & $0.622 \pm 0.017$ \\
Set~3 (medium, baseline)     & $0.631 \pm 0.027$ & $0.618 \pm 0.023$ \\
Set~4 (large)                & $0.618 \pm 0.029$ & $0.620 \pm 0.020$ \\
\midrule
\multicolumn{3}{l}{\textit{Hyperparameters (Set~3)}} \\
Main result                  & 0.631             & 0.618             \\
Learning-rate check          & 0.630             & 0.614             \\
Batch-size check             & 0.641             & 0.619             \\
$\lambda$ check (8 settings) & 0.623--0.632      & 0.612--0.620      \\
\bottomrule
\end{tabular}
\caption{\footnotesize{PINN and PINN-$k$ test $R^2$ across feature sets and hyperparameter checks. Feature-set values report the 5-fold mean $\pm$ SD; hyperparameter values show single evaluation runs or the range across eight physics-weight ($\lambda$) values.}}
\label{tab:pinn-conditions}
\end{table}

% --- Table 3 (was Table 8): plausibility comparison, plain PINN vs. PINN-k. Now the last of
% the 3 tables kept in the main text (see note above on Tables 3-7 moved to appendix). ---
\begin{table}[H]
\centering
\scriptsize
\renewcommand{\arraystretch}{0.9}
\setlength{\tabcolsep}{16pt}
\begin{tabular}{lcc}
\toprule
\textbf{Metric} & \textbf{Plain PINN } $\Delta y_{\max}$ & \textbf{PINN-}$ \bm{k}$ $\Delta y_{\max}$ \\
\midrule
Mean                               & $+2.50$\,m             & $-0.45$\,m \\
SD                                 & $7.27$\,m              & $1.36$\,m \\
Max                                & $+25.79$\,m            & $+4.95$\,m \\
Min                                & $-30.45$\,m            & $-13.88$\,m \\
Implausible ($<5$\,m or $>70$\,m)  & 82 / 58,073 (0.14\%)   & 0 / 58,073 (0\%) \\
\bottomrule
\end{tabular}
\caption{\footnotesize{Plausibility comparison of predicted maximum height deviations ($\Delta y_{\max}$) between plain PINN and PINN-$k$, pooled across all five folds ($N=58{,}073$ plots). PINN-$k$ produces zero implausible heights on the pooled dataset.}}
\label{tab:plausibility-comparison}
\end{table}

\subsection*{Discussion}
\paragraph*{XGBoost is the most accurate model} overall ($R^2 = 0.674$), ahead of the standard DNN ($0.655$), PINN ($0.631$), and PINN-$k$ ($0.618$; Table~\ref{tab:results-q3}). It wins every fold on the 6-survey cohort and nearly every fold on the 4-survey data (trailing in just 1 of 5, by only $0.007$ $R^2$). Because requiring six clean surveys leaves out steep, high-elevation stands (48 vs.\ 296 compartments), the rest of this chapter uses the broader 4-survey cohort. Four possible reasons for the DNN-to-XGBoost gap were tested and ruled out: feature interactions (DNN already captures 80\% of the non-linear gain), training sample size (no wider gap when subsampling), network size (no gain from larger networks), and unusual terrain (XGBoost actually does worse than DNN there, error ratio $1.312$ vs.\ $1.243$). This sets up the key question for Q3: do PINNs' safety and interpretability make up for their lower accuracy?

\paragraph*{Environmental drivers are balanced, not dominated by one feature.}
We checked whether PINN-$k$ leans on one hidden feature the way GNNWR leaned on canopy cover in Q2, using permutation importance (10 repeats per feature; Figure~\ref{fig:terrain-importance}). Here, influence is spread out: the top predictor (wind exposure) accounts for only 15.2\% of total importance, with all eleven terrain and climate variables ranging from 4.7\% to 15.2\% (fold 0). On a typical stand (Plot 119937; Figure~\ref{fig:good-example}), this spread-out adjustment still improves on the regional curve.

% Side by side: both narrow (0.7 width alone), both referenced in the paragraph just above.
\begin{figure}[!htbp]
\centering
\begin{subfigure}[t]{0.58\textwidth}
\centering
\includegraphics[width=\textwidth]{figures/fig_results/q3/q3_pinn_terrain_importance.png}
\caption{}
% BUILT, stable. Source: RESULTS_TABLE.md #6.
\label{fig:terrain-importance}
\end{subfigure}
\hfill
\begin{subfigure}[t]{0.38\textwidth}
\centering
\includegraphics[width=\textwidth]{figures/fig_results/q3/q3_pinn_example_plot_curve.png}
\caption{}
% BUILT 2026-08-24, extended to show plain PINN alongside PINN-k, same treatment as 77226.
% Underlying graph itself also made narrower (5x5, was 7x5) and legend moved to bottom-right,
% same date, see build_good_example_plot_curve() in build_q3_redraft_figures.py.
\label{fig:good-example}
\end{subfigure}
\caption{\small{(a) Share of total permutation importance carried by each terrain/wind feature
in PINN-$k$'s personalised $y_{max}$, fold 0. (b) Plot 119937 (compartment 2040, an ordinary
compartment - rank 91 of 178, yield class 16, close to the population mean): both models'
curves fit the observed points better than the population curve, PINN-$k$ most closely of all.
Not an extreme case, shown because personalisation should work sensibly on a typical plot,
not just the flagged ones.}}
\end{figure}

\paragraph*{A small environmental feature set captures most of the signal}
We tested PINN on four environmental setups with 5-fold cross-validation: no site data, then three growing feature sets (Table~\ref{tab:pinn-conditions}, Figure~\ref{fig:env-sweep}). Even a small set (Set~2) gives an immediate jump over the no-environment baseline, in every fold, for both models. Bigger sets add nothing more, Set~2, Set~3, and Set~4 all score the same within fold-level noise. With site data, PINN-$k$ is more consistent fold to fold than plain PINN; without it, both fall to the same baseline ($R^2 \approx 0.57$). So environmental data matters, but which set barely does; Set~3 is kept as the default mainly for consistency with the rest of the study.

\begin{figure}[!htbp]
\centering
\includegraphics[width=0.75\textwidth]{figures/fig_results/q3/q3_env_sweep_r2.png}
\caption{\small{PINN and PINN-$k$ accuracy against environmental feature-list size.}}
% BUILT, stable. Source: RESULTS_TABLE.md #5.
\label{fig:env-sweep}
\end{figure}

\paragraph*{How PINN shifts curve parameters and matches yield classes}
Compartments were picked using a fixed ranking rule, not by eye (Table~\ref{tab:productivity-top5}). Among 178 compartments with $\geq 50$ pooled test plots, Compartment 1129 had the single largest plot departure ($+25.8$\,m; rank 1/178), and Compartment 2021 had both the largest mean shift and the largest negative departure (Figure~\ref{fig:ymax-deviation-map}). All compartments referenced here are shown on a master map (Appendix Figure~\ref{fig:all-compartments-reference}).

The top three over- and under-productive compartments agree closely between models (Appendix Figure~\ref{fig:hotspot-summary}; Table~\ref{tab:productivity-flagships}). But plain PINN and PINN-$k$ route the adjustment differently (Figure~\ref{fig:ymax-k-deviation-maps}): PINN-$k$ keeps mean $y_{\max}$ shifts small ($\pm 3$\,m) and pushes most of the signal through growth rate ($k$) instead, almost always upward. Four examples show this: Compartment 1045 adjusts mainly through $y_{\max}$, 2057 through $k$, 1070 through a negative $y_{\max}$ shift alone, and 2021 splits the adjustment across both.

\begin{figure}[!htbp]
\centering
\includegraphics[width=0.8\textwidth]{figures/fig_results/q3/q3_pinn_ymax_deviation_map.png}
\caption{\small{Compartment-level asymptote departures under standard PINN (pooled across 5 folds). Mean $y_{\max}$ deviation from the regional curve, highlighting flagship compartments 1129 and 2021.}}
\label{fig:ymax-deviation-map}
\end{figure}

\begin{figure}[!htbp]
\centering
\includegraphics[width=1.1\textwidth]{figures/fig_results/q3/q3_pinn_k_pooled_ymax_k_deviation_maps.png}
\caption{\small{PINN-$k$, pooled across all 5 folds (58,073 plots): mean $y_{max}$ deviation
(left) vs.\ mean $k$ deviation (right), by compartment. $y_{max}$ stays close to flat almost
everywhere ($\pm$3\,m); $k$ shows a structured, almost entirely one-sided pattern instead,
the same kind of signal, routed through a different parameter than plain PINN uses, and one
that barely moves downward at all. Four compartments labelled A-D on both panels (legend below
the figure) to show that different compartments route the signal differently, not just to
repeat the single most extreme case four times.}}
\label{fig:ymax-k-deviation-maps}
\end{figure}

To check biological plausibility, predictions were compared against yield-class records withheld from training. The five most under-productive compartments matched low yield classes in 4 of 5 cases; the five most over-productive matched high yield classes in 3 of 5. Across the full test set, predicted deviations correlate significantly with yield class (Table~\ref{tab:yieldclass-correlation}; Appendix Figure~\ref{fig:yieldclass-scatter}), showing the network picks up broad ecological patterns as well as local terrain effects.

\paragraph*{Why some compartments disagree with yield class}
The over-productive stands that do not match their regional yield class follow a clear pattern (Table~\ref{tab:productivity-top5}): the two mismatched compartments (2057, 1027) are older (60 and 72 years) and wind-exposed, while the three matching ones are younger (25--35 years) and sheltered. Wind exposure is the strongest terrain driver overall (15.2\% permutation importance; Figure~\ref{fig:terrain-importance}), but age alone is not the answer, PINNs are actually more accurate in older stands (Table~\ref{tab:age-height-error}), so the real cause is still unclear. Compartment 2057 also has PINN-$k$'s single largest growth-rate deviation, meaning two independent model setups flag the same stand.

\paragraph*{Biological plausibility differs sharply between PINN variants}
Plain PINN sometimes predicts extreme, biologically impossible heights for a small number of stands (Figure~\ref{fig:implausible-example}). PINN-$k$ never does this, staying realistic across all $58{,}073$ test plots (Table~\ref{tab:plausibility-comparison}, Figure~\ref{fig:plausibility-distribution}).

This happens because plain PINN can only adjust the height ceiling, so a strong environmental signal pushes that ceiling too high to fit fast early growth. PINN-$k$ avoids this by putting that fast growth into the rate parameter ($k$) instead, leaving the ceiling stable. This is not just a fluke of cross-validation: even training directly on the affected compartments (four of five folds) still gave the same inflated ceiling (Table~\ref{tab:compartment1129-crossfold}). Whether these stands are genuinely exceptional or reflect old measurement errors is still an open question, but matching independent yield classes supports real site variation.

\begin{figure}[!htbp]
\centering
\begin{subfigure}[t]{0.48\textwidth}
\centering
\includegraphics[width=\textwidth]{figures/fig_results/q3/q3_pinn_implausible_example_plot.png}
\caption{}
\label{fig:implausible-example}
\end{subfigure}
\hfill
\begin{subfigure}[t]{0.48\textwidth}
\centering
\includegraphics[width=\textwidth]{figures/fig_results/q3/q3_plausibility_distribution.png}
\caption{}
\label{fig:plausibility-distribution}
\end{subfigure}
\caption{\small{(a) Plot 77226 (compartment 1129): observed heights against the population
curve, plain PINN's curve, and PINN-$k$'s curve. Plain PINN pushes the ceiling to an implausible
77.8m; PINN-$k$ keeps the ceiling modest (53.0m) but personalises growth rate instead
($k$=0.0222, the dataset's single highest value), the two curves cross around age 70,
PINN-$k$ fitting the young observed points better before plain PINN's ceiling pulls ahead.
(b) Distribution of $y_{max}$ deviation from the population curve, plain PINN vs.\ PINN-$k$,
pooled across all 5 folds (58,073 plots each). Same underlying signal, very different spread.}}
\end{figure}

\paragraph*{Model performance differs at age and height extremes}
Accuracy shifts across a stand's life (Table~\ref{tab:age-height-error}). In young stands ($0$--$10$\,m), the unconstrained DNN wins (MAE $2.68$\,m vs.\ $3.37$\,m for plain PINN, $3.15$\,m for PINN-$k$), since PINN's fixed curve shape struggles with early growth. In older stands ($>75$ years, $n=751$), it flips: the DNN is worst ($6.83$\,m vs.\ $6.40$\,m for plain PINN), since nothing stops it drifting without a physical curve to anchor it. This matters for timber value: volume scales with height to the power $1.48$--$1.69$, so a $10$\% height error becomes a $15$--$17$\% volume error. Since PINN's errors sit mostly in young stands, its volume predictions stay reliable where commercial value is highest.

\paragraph*{Balancing predictive accuracy against biological safety}
Comparing plain PINN (height only) with PINN-$k$ (height and growth rate) shows a clear trade-off: plain PINN is slightly more accurate in every setting tested (Table~\ref{tab:pinn-conditions}), but PINN-$k$ models both growth parameters and never predicts an impossible height (Table~\ref{tab:plausibility-comparison}). Still, the two parameters move together sensibly ($r = 0.43$): positive growth-rate shifts pair with negative height shifts across compartments (Figure~\ref{fig:ymax-k-deviation-maps}). Hyperparameter checks confirm plain PINN's edge holds across different settings (Table~\ref{tab:pinn-conditions}). Turning the physics loss off entirely ($\lambda = 0$) barely changes accuracy (PINN-$k$ $R^2 = 0.620$ vs.\ $0.618$ at the default $\lambda = 1$), with a small drop once $\lambda > 1$, physics guidance shapes \emph{how} the model fits, not how well.

\subsection*{Q3: Conclusion}
XGBoost is the most accurate model for top height, but physics-guided networks trade modest accuracy for safer predictions: PINN-$k$ eliminates biologically impossible heights entirely, while plain PINN produces rare but extreme errors (82 of 58,073 plots, 0.14\%; Table~\ref{tab:plausibility-comparison}). Unlike Q2, Q3's environmental drivers are balanced across terrain and climate features, and performance varies by stand age: unconstrained DNNs lead in young stands, but PINN's fixed ceiling is more reliable in mature ones. As shown above, physics guidance does not boost raw accuracy, but it enforces biological plausibility ($99.86$\% of plots for plain PINN, $100$\% for PINN-$k$) at a minor cost to precision.

%=======================================================================================
%=======================================================================================

%=======================================================================================
%=======================================================================================
\chapter{Conclusion and Future Work}
\label{ch:conc}

\section{Main conclusions}

Canopy cover was the strongest correlate of departure from the shared growth curve (Q1). Management variables added smaller signals. Among the environmental variables, only $\mathit{topex}$ and wind exposure showed repeatable local associations after compartment-level differences were considered. These are observational associations, not causal effects. Canopy cover should therefore not be interpreted as a proven management lever.

For the related $\Delta y_{\max}$ target, GNNWR provided only a small improvement over the global models (Q2). Canopy cover accounted for approximately 80--90\% of its explained signal, based on the reduction in $R^2$ when it was removed. The added spatial complexity therefore revealed little further environmental structure.

XGBoost produced the most accurate top-height predictions (Q3). PINN-$k$ had an RMSE approximately 0.4\,m higher, but produced no predictions outside the predefined 5--70\,m plausibility range. Plain PINN produced a small number of large violations. Physics guidance therefore improved plausibility rather than predictive accuracy. This finding applies only within the ages and forest conditions evaluated.

Together, the results address one part of the Forest Digital Twin concept: combining repeated observations and models to explain and predict forest change \cite{baskent2026Review_C1_SS,buonocore2022Proposal_C1_SS}. Q1 and Q2 identify conditions associated with unusual growth. Q3 tests whether those signals support reliable height prediction. Real-time updating, disturbance simulation, management scenarios, and species beyond Sitka spruce remain outside the scope of this work.

\section{Limitations}

\paragraph*{Data and targets}

The environmental layers were coarse relative to individual plots. SoilGrids has a 250\,m resolution, while CHELSA and HadUK climate data are approximately 1\,km \cite{poggio2021SoilGrids_C3_LID,karger2021CHELSA_C3_LID,hollis2019HadUKGrid_C3_LID}. These layers cannot represent fine-scale soil variation or short events such as storms. This may explain the weak and unstable soil and climate signals.

Historical uncertainty also affects the compartment records. Current species, yield-class, and thinning records were applied to earlier surveys and may not describe each compartment accurately at every date. Compartments 2057 and 1027, already flagged in Q3 for disagreeing with their assigned yield class, are exactly the kind of case this could produce, though this project cannot confirm whether their yield-class labels have actually changed over time. Q2's $\Delta y_{\max}$ target has a further limitation: with only four to six observations per plot, differences in growth rate and final height cannot be separated reliably. It is therefore a summary of trajectory departure, not a direct measurement of maximum height.

The six-survey cohort contains only 47 compartments, compared with 231 in the main cohort. Its weak performance may reflect limited spatial variation and unstable holdout folds rather than a fundamental model failure.

\paragraph*{Models and evaluation}

Q1 used one five-fold spatial partition. Its findings, including the slope/$\mathit{topex}$ stability contrast at the centre of Q1's Elastic Net/LMM disagreement, may therefore depend partly on the selected split. GNNWR was tested with additional partitions, but training randomness could not be separated from changes in fold membership, so its reported six-survey instability cannot be cleanly attributed to the model itself. Spatial blocking reduces leakage between nearby observations, but reliable evaluation still depends on having enough independent spatial groups \cite{roberts2017Cross_C5_EVAL,ploton2020Spatial_C5_EVAL}.

Elastic Net and the mixed-effects model assume linear relationships unless nonlinear terms are supplied. This assumption does not fully match the curved responses observed for slope and wind exposure, which qualifies Q1's finding that slope's coefficient sign agrees across models, agreement on sign does not rule out both models missing the same curvature, though this is likely a minor effect here since canopy cover dominates the explained variance. GNNWR was not systematically tuned, so its small advantage over the global models in Q2 remains uncertain. PINN and PINN-$k$ were also evaluated on one cohort and one training run. The lambda ablation behind Q3's headline claim that physics guidance shapes fit rather than accuracy has not been checked on the six-survey cohort or across seeds, and the separate contribution of the Chapman--Richards branch was not isolated.

\section{Future work}

The first priority is stronger evaluation. Future experiments should repeat all models across several seeds, tune GNNWR and PINN-$k$ (each GNNWR run took roughly two hours, so cheaper equivalent models with a similar locally-varying structure are also worth testing), and test whether the Chapman--Richards branch contributes beyond the unconstrained neural component. Training on earlier surveys and testing on later surveys would measure temporal prediction, most naturally on the six-survey cohort, since it has the most years per plot, though that is also the cohort with too few compartments (48/296) for reliable spatial evaluation, so a temporal split would inherit the same small-sample problem, while evaluation in another forest, and in other commercially important species beyond Sitka spruce, would test geographic and species transfer.

Future attribution should also examine different forms of growth departure. Separate targets could represent persistently high or low growth, early acceleration, later plateauing, disturbance, and recovery. More surveys would help distinguish changes in growth rate from changes in height potential.

A fuller Forest Digital Twin could move from plot summaries to individual trees. Multi-temporal LiDAR can track tree-level growth and loss, although matching trees between surveys remains difficult \cite{duncanson2018Utility_FW}. Using complete three-dimensional point clouds could represent crown size, competition, structural damage, and mortality rather than top height alone \cite{weiser2022Individual_FW}.

Finally, finer soil and climate measurements could be combined with satellite embeddings such as AlphaEarth. These embeddings integrate spatial, temporal, and remote-sensing information, but would reduce direct interpretability and still require spatial and temporal validation \cite{brown2025AlphaEarth_C4_NN}.

%===========================


%=======================================================================================
%=======================================================================================
\bibliographystyle{plain}
\bibliography{mybibfile}
%=======================================================================================

\appendix

% =============================================================================
% Working extract, full dissertation appendix (Data, Methodology, and Results,
% all three Qs). Renamed from appendix_results.tex 2026-08-23 and scope widened,
% the old Results-only restriction below no longer applies, since Data/Methodology
% appendix material is now being built up here too, not deferred.
% Content fragment, NOT a standalone compilable document, meant to be merged
% into the main file's existing \chapter{Appendix} \label{chap:appendix} in
% 20th_1_31_pm_audited.tex (currently an empty placeholder chapter, see its own
% % AUDIT: note listing 7 unresolved \ref calls).
%
% Structure: \chapter{Data}, \chapter{Methodology}, then one \section per research
% question for Results-appendix material (matching the three working results_qN_*.tex
% files). Built content is filled in as it's ready; not-yet-built content is a short
% placeholder note, not invented numbers.
% =============================================================================

\chapter{Appendix: Data}

\paragraph*{Field derivations and forestry definitions}
{\scriptsize
\setlength{\tabcolsep}{4pt}
\renewcommand{\arraystretch}{1.15}
\begin{longtable}{p{2.4cm}p{11.4cm}}
\caption{Full derivation and forestry definitions for fields referenced in Table~\ref{tab:data-fields} but not spelled out there, one field per block. Units given wherever the source defines them.}
\label{tab:appendix-formulas} \\
\noalign{\hrule height 1pt}
\endfirsthead
\multicolumn{2}{l}{\normalfont\itshape (continued from previous page)} \\
\noalign{\hrule height 1pt}
\endhead
\noalign{\hrule height 1pt}
\multicolumn{2}{r}{\normalfont\itshape (continued on next page)} \\
\endfoot
\noalign{\hrule height 1pt}
\endlastfoot
\multicolumn{2}{l}{\textbf{\texttt{Age}}} \\
Formula & $\texttt{LiDAR\_year} - \texttt{plyr}$ \\
Meaning & Stand age at the survey date, in years since planting \\
Used in project & Modelling input; cleaning filter. \fc{Field-description doc states valid range 0--150 years; \texttt{clean\_master\_data.py} implements 0--200, documentation/code mismatch, unresolved} \\
\noalign{\hrule height 1pt}
\multicolumn{2}{l}{\textbf{\texttt{Top\_Height95}}} \\
Formula & $\texttt{elev\_percentile\_95th}\times1.1$ \\
Meaning & FR's corrected estimate of Top Height (m): mean height of the 100 largest-diameter trees per hectare. The $\times1.1$ factor is FR's stated correction for known underestimation at the 95th return percentile \\
Used in project & Audit ingredient only (feeds \texttt{Vol95}/\texttt{GYCspec95}). This project's target is the \emph{uncorrected} \texttt{elev\_percentile\_95th}, a deliberate deviation from FR's own corrected convention \\
\noalign{\hrule height 1pt}
\multicolumn{2}{l}{\textbf{\texttt{Vol95}}} \\
Formula & $(g_1\cdot\texttt{Top\_Height95}^{g_3})\times(\texttt{area}/10000)\times\texttt{CanopyCover}$ \\
Meaning & Stand volume (m\textsuperscript{3}), scaled from a per-ha allometric relationship by plot area (ha) and canopy cover fraction \\
Used in project & Used downstream as a disturbance-screening signal. $g_1$/$g_3$ are pre-assigned per plot in the source FR export, not computed or selected by this project, e.g.\ $g_1{=}3.509144,\,g_3{=}1.478136$ and $g_1{=}2.922183,\,g_3{=}1.694732$ are two real examples. % Verified 2026-08-23 against clean_master_4survey.parquet: SS alone has 10 distinct (g1,g3) pairs, not split cleanly by thinning status (6 of the 10 pairs occur under BOTH Thin=0 and Thin=1) or by yield class (each pair spans 7-12 different yldc values), the true determinant isn't visible in this project's own tables, likely region or an internal FR yield-model version. Do not describe this as "N pairs per thinning status" or imply thinning/yldc determines the pair. \\
\noalign{\hrule height 1pt}
\multicolumn{2}{l}{\textbf{\texttt{Vol\_RM95}}} \\
Formula & $(-35.8733+5.3486\cdot\texttt{Top\_Height95}^{1.5424})\times(\texttt{area}/10000)\times\texttt{CanopyCover}$ \\
Meaning & Robertson \& Miller generic volume equation (m\textsuperscript{3}), species-independent alternative to \texttt{Vol95} \\
Used in project & Not kept in the master export; not used \\
\noalign{\hrule height 1pt}
\multicolumn{2}{l}{\textbf{\texttt{GYCspec95}}} \\
Formula & $\bigl(\texttt{Top\_Height95}/(1-e^{-p_4\cdot\texttt{Age}})^{p_5}-p_1-p_3\times2\bigr)/p_2$ \\
Meaning & General Yield Class (continuous scale; FR's own tables report GYC in multiples of 2, e.g.\ YC14, YC16): potential productivity, the maximum mean annual increment of cumulative volume (m\textsuperscript{3}/ha/yr) for the species on that site under standard management \\
Used in project & Never used downstream; would be leakage if it were (contains \texttt{Top\_Height95}) \\
\noalign{\hrule height 1pt}
\multicolumn{2}{l}{\textbf{$y_{\max}^{\texttt{yldc}}$ (RQ3)}} \\
Formula & $\texttt{yldc}\cdot p_2+p_1+p_3\times2$ \\
Meaning & The height ceiling (m) that the plot's independently assigned yield class implies, under the same site-index model as \texttt{GYCspec95}, solved the opposite direction \\
Used in project & Used, RQ3's independent comparison baseline (\texttt{data\_processing/export\_growth\_curve\_tables.py}) \\
\noalign{\hrule height 1pt}
\multicolumn{2}{l}{\textbf{\texttt{DBHmean}} \ct{source PDF/report, not yet in citation list}} \\
Formula & $(b_1+b_2\cdot\texttt{Spacing})\times\texttt{TopHeight}^{b_4}$ \\
Meaning & Mean diameter at breast height, 1.3\,m above ground (cm); needed to estimate stems per hectare \\
Used in project & \textbf{Not present anywhere in this project's data or code.} Supporting context only, for the claim that top height feeds an established forestry allometric toolkit (SS, thinned: $b_1{=}0.714157,\,b_2{=}0.239399,\,b_4{=}1.079545$; unthinned: $b_1{=}1.212457,\,b_2{=}0.609335,\,b_4{=}0.754978$) \\
\noalign{\hrule height 1pt}
\multicolumn{2}{l}{\textbf{\texttt{BA}} \ct{same source}} \\
Formula & $h_1\times\texttt{TopHeight}^{h_4}$ \\
Meaning & Basal area (m\textsuperscript{2}/ha): cross-sectional stem area summed per hectare \\
Used in project & Not present in this project's data. Same supporting-context caveat (SS, thinned: $h_1{=}7.499437,\,h_4{=}0.497925$; unthinned: $h_1{=}5.517392,\,h_4{=}0.771801$) \\
\noalign{\hrule height 1pt}
\multicolumn{2}{l}{\textbf{Volume (source PDF, Eq.~13)} \ct{same source}} \\
Formula & $g_1\times\texttt{TopHeight}^{g_3}$ (no area/canopy scaling) \\
Meaning & Volume over bark, per hectare (m\textsuperscript{3}/ha) \\
Used in project & \fc{Distinct equation from this project's own \texttt{Vol95} above, same $g_1$/$g_3$ notation, different formula (no area/\texttt{CanopyCover} scaling). Do not conflate the two in the appendix text} \\
\noalign{\hrule height 1pt}
\multicolumn{2}{l}{\textbf{Top Height (forestry definition)}} \\
Formula &, \\
Meaning & Mean height of the 100 largest-diameter trees per hectare (m) \\
Used in project & \texttt{elev\_percentile\_95th} is an ALS-derived proxy for this classical mensuration concept, not identical to it by construction \\
\noalign{\hrule height 1pt}
\multicolumn{2}{l}{\textbf{Yield Class (forestry definition)}} \\
Formula &, \\
Meaning & Index of site productivity: max mean annual increment of cumulative volume (m\textsuperscript{3}/ha/yr); tabulated in multiples of 2 in FR yield tables, continuous here \\
Used in project & \texttt{yldc} is this project's realisation of it \\
\end{longtable}
}

\paragraph*{Environmental variable definitions and coverage}
% Filed under Data, not Methodology, despite the source draft's "METHODLOGOY" heading,
% this table's own caption says feature-selection/VIF/final model membership are described
% in the Methodology chapter, "not here". What's actually here (extraction method, coverage,
% distinct values, min/max per raw variable) is Data-chapter material, variable definitions,
% not model construction. Moved 2026-08-23; flag if this placement is wrong.
{\tiny
\setlength{\tabcolsep}{2pt}
\renewcommand{\arraystretch}{0.9}
\begin{longtable}{p{2.3cm}p{1.6cm}p{3.0cm}p{1.3cm}p{1.1cm}p{2.3cm}}
\caption{Per-variable definitions, extraction method and coverage for the environmental table
(71,766 plots). Source: \texttt{methodlogy\_env\_setpick.md}, verified 2026-08-10 against
\texttt{plot\_environmental\_features.parquet}. Feature-selection, VIF screening and which
variables entered a final model set are described in Chapter~\ref{ch:methodology}, not here.}
\label{tab:data-env-variables_detail_app} \\
\hline
Variable & Family & Extraction method & Valid/ Missing & Distinct & Min--Max \\
\hline
\endfirsthead
\multicolumn{6}{l}{\normalfont\itshape (continued from previous page)} \\
\hline
Variable & Family & Extraction method & Valid/ Missing & Distinct & Min--Max \\
\hline
\endhead
\hline
\multicolumn{6}{r}{\normalfont\itshape (continued on next page)} \\
\endfoot
\hline
\endlastfoot
\texttt{elevation} & Terrain & Direct raster read & 71{,}766/0 & 4{,}028 & 18.2--561.1\,m \\
\texttt{slope\_degrees} & Terrain & $\arctan(\sqrt{dz_{dx}^2+dz_{dy}^2})$, \texttt{np.gradient} & 71{,}766/0 & 15{,}326 & 0.0--42.84$^\circ$ \\
\texttt{northness}, \texttt{eastness} & Terrain & $\cos/\sin$(aspect, radians) & 71{,}766/0 & 15{,}387/15{,}572 & $-1$--$1$ \\
\texttt{profile\_/plan\_curvature} & Terrain & 2nd derivative of DTM \citep{zevenbergen1987} & 71{,}766/0 & 16{,}835/16{,}829 & $-0.007$--$0.008$ / $-0.67$--$0.40$ \\
\texttt{local\_relief\_500m} & Terrain & Max$-$min elevation, 21$\times$21-cell window & 71{,}766/0 & 6{,}086 & 7.5--461.9\,m \\
\texttt{solar\_radiation\_index} & Terrain & $\cos(\text{slope})\cos(\text{zenith})+\sin(\text{slope})\sin(\text{zenith})\cos(\text{aspect}-\text{azimuth})$, no horizon shading & 71{,}766/0 & 16{,}012 & 0.338--1.000 \\
\texttt{frost\_hollow\_flag} & Terrain & Binary: TPI $<$ own 15th pctile AND plan curvature $<0$ & 71{,}766/0 & 2 & 0 ($n{=}62{,}691$)/1 ($n{=}9{,}075$) \\
\texttt{ceh\_twi} & Terrain & Direct raster read (CEH TWI) & 71{,}727/39 & 16{,}800 & 3.75--16.86 \\
\texttt{topex} & Wind & Sum of horizon angles, 8 directions \citep{wilson1984} & 71{,}766/0 & 16{,}839 & $-85.9$--$68.6^\circ$ \\
\texttt{windward\_topex} & Wind & Same, single 225$^\circ$ bearing & 71{,}766/0 & 8{,}093 & $-19.8$--$19.4^\circ$ \\
\texttt{gwa\_weibull\_a/k\_10m/50m} & Wind & Direct raster read, national grid & 71{,}766/0 & 1{,}816--1{,}824 & varies by variable \\
\texttt{gwa\_wind\_speed\_10m} & Wind & Nearest-cell extraction, REST API & 71{,}766/0 & 1{,}816 & 0.069--11.18\,m/s \\
\texttt{whcl} & Wind & Direct field read (raw GPKG) & 71{,}766/0 & 6 of 7 possible & Class 1 absent \\
\texttt{soilgrids\_ph} & Soil & Direct read, topsoil 0--5\,cm & 71{,}353/413 & 13 & 4.0--5.2 \\
\texttt{ceh\_pedotope} & Soil & Categorical, one-hot & 71{,}727/39 & 7 levels &, \\
\texttt{ceh\_subsurface\_drainage} & Soil & Categorical, one-hot & 71{,}727/39 & 3 levels &, \\
\texttt{ceh\_textural\_composition} & Soil & Categorical, one-hot & 71{,}727/39 & 4 levels &, \\
\texttt{dist\_to\_watercourse} & Soil/spatial & Distance to nearest OS Open Rivers line & 71{,}766/0 & 71{,}758 & 0.0045--1850.9\,m \\
\texttt{chelsa\_gdd5\_degc} & Climate & Direct read, 1981--2010 & 71{,}766/0 & 206 & 906.5--1580.8 \\
\texttt{chelsa\_bio12\_precip\_mm} & Climate & Direct read, 1981--2010 & 71{,}766/0 & 269 & 1729.8--3004.8 \\
\texttt{tas\_mean}, \texttt{groundfrost\_mean} & Climate & Cohort-suffixed, resolved to cohort's real years & 71{,}766/0 & 149 each & 5.75--9.15 / 8.76--9.47 \\
\texttt{dist\_to\_scpt\_/\_block\_boundary} & Spatial & Distance to own sub-cpmt/block polygon & 71{,}766/0 & 68{,}188/53{,}974 & $\sim$0--217\,m \\
\texttt{dist\_to\_forest\_perimeter} & Spatial & Distance to buffered union of compartments & 71{,}766/0 & 71{,}701 & $\sim$0--764.7\,m \\
\texttt{cpmt\_compactness\_ratio} & Spatial & Perimeter/area, per compartment & 71{,}766/0 & 296 & 0.0097--0.240 \\
\texttt{dist\_to\_road} & Spatial & Distance to nearest OS Open Roads line & 71{,}766/0 & 71{,}580 & 0.056--2375.7\,m \\
\texttt{CanopyCover} & Stand & Mean fraction across survey years & 71{,}766/0 & 70{,}769 & 0.008--0.996 \\
\texttt{time\_since\_thinning} & Stand & NaN-filled with 0 downstream & 71{,}766/0 & 19 & 0--16\,yr \\
\end{longtable}
}

\begin{table}[H]
\centering
\textbf{External data source details}\label{app:data-sources}\\[4pt]
{\small Table~\ref{tab:data-env-sources} in the Data chapter summarises the seven data families used; this table gives the access method, licensing, and temporal vintage behind each one.\\[6pt]}
\scriptsize
\setlength{\tabcolsep}{3pt}
\renewcommand{\arraystretch}{1.1}
\begin{tabular}{p{1.8cm}p{4.7cm}p{2.0cm}p{5.2cm}}
\hline
Data family & Provider and access & Licence & Temporal vintage and caveats \\
\hline
Terrain & OS Terrain 50 DTM, direct raster read & OS OpenData & Single static snapshot; not matched to individual survey years \\
Soil & CEH Natural Capital pedotope raster (one of 9 layers in a single downloaded zip); SoilGrids v2.0 (ISRIC, public Cloud-Optimised GeoTIFF via \texttt{rasterio}'s \texttt{/vsicurl/}, no authentication) & CEH: institutional download; SoilGrids: public, no auth & Static; SoilGrids pH uses only the shallowest of six depth bands (0--5\,cm topsoil) \\
Wind exposure & Global Wind Atlas GBR national rasters (Weibull-A/K, mean wind speed), direct raster read and public REST API & Public & 2006--2015 wind climatology: one fixed value per plot location, not recalculated per survey year \\
Observed wind & MIDAS Open weather stations & Public (Met Office) & Computed per survey-to-survey interval from hourly station observations, not a static layer \\
Long-term climate & CHELSA v2.1, direct public download, no authentication & Public & 1981--2010 climatology (terrain-conditioned downscaling, not plain interpolation); static, not tied to any survey year \\
Matched climate & HadUK-Grid, CEDA & Public & Genuinely multi-year: averaged across each cohort's own real survey years (2008/2012/2021/2023 for the 4-survey cohort; all six years for the 6-survey cohort), not a single reused snapshot \\
Spatial context & OS Open Roads and Open Rivers, vector line data & OS OpenData & No stated vintage in this project's own downloaded package docs; the raw file's download timestamp is $\sim$July 2026, applied retroactively to every survey year back to 2002. Minor for rivers (stable over time); more relevant for roads, since forestry access tracks can be built within that timeframe \\
\hline
\end{tabular}
\caption{External data source access method, licensing, and temporal vintage, one row per
data family in Table~\ref{tab:data-env-sources}.}
\label{tab:data-sources-detail}
\end{table}


\chapter{Appendix: Methodology}

\paragraph*{Feature-set membership by research question}
\label{app:feature-sets}
% Converted from raw notes 2026-08-23. RSQ->Q mapping verified against actual production
% code (not just target-name matching): RSQ1 loads via models/common/torch_data.py's
% ENV_TERRAIN_FEATURE_SETS, which feeds Q3's DNN/PINN models directly (target
% elev_percentile_95th, confirmed both by that file's own TARGET_COLUMN and by
% feature_set_builder.py's own docstring: "RSQ1: ...feeds dnn_env_terrain/pinn_env_terrain*").
% RSQ2 loads in run_rq2_attribution.py, Q1's real production script (target
% mean_cr_residual). RSQ3 loads in run_rq3_gnnwr.py, Q2's real GNNWR production script
% (target local_y_max_difference). No mismatches found across every script checked.
%
% Machine-readable source: documentation/env_feature_sets_manifest.csv (225 rows), the
% snapshot below is for readability; the CSV is the definitive record if the two disagree.

Ranking and screening is repeated separately for each question's own target, so final Set
membership differs even where two questions share a Set number. Sets are not strictly nested:
VIF-based collinearity screening (variables with variance inflation factor $>5$ removed
iteratively) is re-run each time the candidate pool expands, so a variable present in a smaller
Set can be screened out once a larger Set's wider pool makes it redundant. Elevation is the
clearest example: it survives Set2 and Set3 in every question below, but is removed from Set4
once the wider terrain/wind pool is considered together, this reflects statistical redundancy
with other terrain/wind variables (elevation is a well-known confound for temperature, wind
exposure and radiation), not that elevation is ecologically unimportant.

\begin{table}[H]
\centering
\textbf{Q1 (target \texttt{mean\_cr\_residual})}\\[4pt]
\scriptsize
\renewcommand{\arraystretch}{1.1}
\begin{tabularx}{\linewidth}{@{} l r >{\raggedright\arraybackslash}X @{}}
\toprule
Set & $n$ & Columns \\
\midrule
Set1 & 4 & \texttt{CanopyCover}, \texttt{time\_since\_thinning}, \texttt{time\_since\_thinning\_missing}, \texttt{recent\_thinning\_5yr} \\
Set2 & 14 & Set1 + \texttt{chelsa\_bio12\_precip\_mm}, \texttt{chelsa\_gdd5\_degc}, \texttt{gwa\_weibull\_k\_50m}, \texttt{slope\_degrees}, \texttt{gwa\_weibull\_a\_10m}, \texttt{local\_relief\_500m}, \texttt{eastness}, \texttt{dist\_to\_road}, \texttt{gwa\_weibull\_a\_50m}, \texttt{tas\_mean} \\
Set3 & 14 & Set1 + \texttt{gwa\_weibull\_k\_50m}, \texttt{slope\_degrees}, \texttt{gwa\_weibull\_a\_10m}, \texttt{local\_relief\_500m}, \texttt{eastness}, \texttt{gwa\_weibull\_a\_50m}, \texttt{ceh\_twi}, \texttt{whcl}, \texttt{topex}, \texttt{gwa\_wind\_speed\_10m} \\
Set4 & 19 & Set3 + \texttt{chelsa\_bio12\_precip\_mm}, \texttt{dist\_to\_road}, \texttt{tas\_mean}, \texttt{dist\_to\_scpt\_boundary}, \texttt{soilgrids\_ph} \\
Set5 & 28 & Set1 + \texttt{slope\_degrees}, \texttt{northness}, \texttt{eastness}, \texttt{profile\_curvature}, \texttt{plan\_curvature}, \texttt{tpi\_500m}, \texttt{local\_relief\_500m}, \texttt{frost\_hollow\_flag}, \texttt{windward\_topex}, \texttt{gwa\_weibull\_a\_10m}, \texttt{gwa\_weibull\_k\_10m}, \texttt{gwa\_weibull\_a\_50m}, \texttt{gwa\_weibull\_k\_50m}, \texttt{dist\_to\_forest\_perimeter}, \texttt{dist\_to\_block\_boundary}, \texttt{cpmt\_compactness\_ratio}, \texttt{dist\_to\_road}, \texttt{dist\_to\_watercourse}, \texttt{gwa\_wind\_speed\_10m}, \texttt{soilgrids\_ph}, \texttt{ceh\_twi}, \texttt{tas\_mean}, \texttt{groundfrost\_mean}, \texttt{whcl} \\
\bottomrule
\end{tabularx}
\caption{\small{Q1 (RSQ2) feature-set membership, elastic\_net/XGBoost/LMM attribution pipeline.}}
\label{tab:feature-sets-q1}
\end{table}

% ceh_* columns are one-hot dummy levels, reference level already dropped.
\begin{table}[H]
\centering
\textbf{Q2 (target \texttt{local\_y\_max\_difference})}\\[4pt]
\scriptsize
\renewcommand{\arraystretch}{1.1}
\begin{tabularx}{\linewidth}{@{} l r >{\raggedright\arraybackslash}X @{}}
\toprule
Set & $n$ & Columns \\
\midrule
Set1 & 4 & \texttt{CanopyCover}, \texttt{time\_since\_thinning}, \texttt{time\_since\_thinning\_missing}, \texttt{recent\_thinning\_5yr} \\
Set2 & 14 & Set1 + \texttt{windward\_topex}, \texttt{elevation}, \texttt{cpmt\_compactness\_ratio}, \texttt{gwa\_wind\_speed\_50m}, \texttt{slope\_degrees}, \texttt{tpi\_500m}, \texttt{dist\_to\_road}, \texttt{topex}, \texttt{solar\_radiation\_index}, \texttt{soilgrids\_ph} \\
Set3 & 12 & Set1 + \texttt{windward\_topex}, \texttt{elevation}, \texttt{gwa\_wind\_speed\_50m}, \texttt{slope\_degrees}, \texttt{tpi\_500m}, \texttt{topex}, \texttt{solar\_radiation\_index}, \texttt{northness} \\
Set4 & 20 & Set1 + \texttt{windward\_topex}, \texttt{gwa\_wind\_speed\_50m}, \texttt{slope\_degrees}, \texttt{tpi\_500m}, \texttt{topex}, \texttt{northness}, \texttt{cpmt\_compactness\_ratio}, \texttt{dist\_to\_road}, \texttt{chelsa\_bio12\_precip\_mm}, \texttt{tas\_mean}, \texttt{soilgrids\_ph}, \texttt{ceh\_subsurface\_drainage=2.0}, \texttt{ceh\_pedotope=2.0}, \texttt{ceh\_pedotope=8.0}, \texttt{ceh\_textural\_composition=2.0}, \texttt{dist\_to\_forest\_perimeter} \\
Set5 & 31 & Set1 + \texttt{slope\_degrees}, \texttt{northness}, \texttt{eastness}, \texttt{profile\_curvature}, \texttt{plan\_curvature}, \texttt{ceh\_twi}, \texttt{frost\_hollow\_flag}, \texttt{windward\_topex}, \texttt{whcl}, \texttt{tpi\_500m}, \texttt{local\_relief\_500m}, \texttt{tas\_mean}, \texttt{groundfrost\_mean}, \texttt{soilgrids\_ph}, \texttt{dist\_to\_watercourse}, \texttt{dist\_to\_forest\_perimeter}, \texttt{dist\_to\_scpt\_boundary}, \texttt{cpmt\_compactness\_ratio}, \texttt{dist\_to\_road}, \texttt{ceh\_pedotope=11.0}, \texttt{ceh\_pedotope=2.0}, \texttt{ceh\_pedotope=5.0}, \texttt{ceh\_pedotope=8.0}, \texttt{ceh\_pedotope=9.0}, \texttt{ceh\_subsurface\_drainage=2.0}, \texttt{ceh\_subsurface\_drainage=3.0}, \texttt{ceh\_textural\_composition=5.0} \\
\bottomrule
\end{tabularx}
\caption{\small{Q2 (RSQ3) feature-set membership, GNNWR spatial-attribution pipeline. Note Set3's non-monotonic size (12 columns, smaller than Set2's 14) is expected, not an error: Set2 and Set3 are built by two different rules (top-10 by rank vs.\ top-half-of-category by rank) over overlapping but not identical candidate pools, Set3 is not defined as a superset of Set2.}}
\label{tab:feature-sets-q2}
\end{table}

\begin{table}[H]
\centering
\textbf{Q3 (target \texttt{elev\_percentile\_95th})}\\[4pt]
\scriptsize
\renewcommand{\arraystretch}{1.1}
\begin{tabularx}{\linewidth}{@{} l r >{\raggedright\arraybackslash}X @{}}
\toprule
Set & $n$ & Columns \\
\midrule
Set1 & 4 & \texttt{CanopyCover}, \texttt{time\_since\_thinning}, \texttt{time\_since\_thinning\_missing}, \texttt{recent\_thinning\_5yr} \\
Set2 & 14 & Set1 + \texttt{gwa\_weibull\_k\_50m}, \texttt{dist\_to\_scpt\_boundary}, \texttt{elevation}, \texttt{tas\_mean}, \texttt{eastness}, \texttt{windward\_topex}, \texttt{gwa\_weibull\_a\_10m}, \texttt{gwa\_weibull\_a\_50m}, \texttt{slope\_degrees}, \texttt{cpmt\_compactness\_ratio} \\
Set3 & 15 & Set1 + \texttt{gwa\_weibull\_k\_50m}, \texttt{elevation}, \texttt{eastness}, \texttt{windward\_topex}, \texttt{gwa\_weibull\_a\_10m}, \texttt{gwa\_weibull\_a\_50m}, \texttt{slope\_degrees}, \texttt{gwa\_weibull\_k\_10m}, \texttt{solar\_radiation\_index}, \texttt{ceh\_twi}, \texttt{gwa\_wind\_speed\_10m} \\
Set4 & 18 & Set1 + \texttt{gwa\_weibull\_k\_50m}, \texttt{eastness}, \texttt{windward\_topex}, \texttt{gwa\_weibull\_a\_10m}, \texttt{gwa\_weibull\_a\_50m}, \texttt{slope\_degrees}, \texttt{gwa\_weibull\_k\_10m}, \texttt{solar\_radiation\_index}, \texttt{ceh\_twi}, \texttt{gwa\_wind\_speed\_10m}, \texttt{dist\_to\_scpt\_boundary}, \texttt{tas\_mean}, \texttt{chelsa\_gdd5\_degc}, \texttt{cpmt\_compactness\_ratio} \\
Set5 & 30 & Set1 + \texttt{slope\_degrees}, \texttt{northness}, \texttt{eastness}, \texttt{profile\_curvature}, \texttt{plan\_curvature}, \texttt{local\_relief\_500m}, \texttt{frost\_hollow\_flag}, \texttt{topex}, \texttt{windward\_topex}, \texttt{gwa\_weibull\_a\_10m}, \texttt{gwa\_weibull\_k\_10m}, \texttt{gwa\_weibull\_a\_50m}, \texttt{gwa\_weibull\_k\_50m}, \texttt{dist\_to\_forest\_perimeter}, \texttt{dist\_to\_scpt\_boundary}, \texttt{cpmt\_compactness\_ratio}, \texttt{dist\_to\_road}, \texttt{dist\_to\_watercourse}, \texttt{gwa\_wind\_speed\_10m}, \texttt{soilgrids\_ph}, \texttt{ceh\_twi}, \texttt{chelsa\_gdd5\_degc}, \texttt{chelsa\_bio12\_precip\_mm}, \texttt{tas\_mean}, \texttt{groundfrost\_mean}, \texttt{whcl} \\
\bottomrule
\end{tabularx}
\caption{\small{Q3 (RSQ1) feature-set membership, DNN/PINN/PINN-$k$ prediction pipeline.}}
\label{tab:feature-sets-q3}
\end{table}

% Renamed RSQ1->Q3, RSQ2->Q1, RSQ3->Q2 per the verified mapping above. "NLME" kept as the
% row label here since that's the actual module/function name in code (models/spatial_
% attribution/nlme.py), the main Methodology chapter calls the same model "linear mixed-
% effects model (LMM)" in prose; both names refer to the identical model, not two different
% ones. SCOPE_GROUPS claim ("no Q1 equivalent") verified by grep: SCOPE_GROUPS only appears
% in models/growth_curve_attribution/*.py (Q2's own module), nothing in models/spatial_
% attribution/ (Q1's module), but Q1 DOES have its own separate category-level ablation
% (the "soil +0.022, climate -0.056" result already in Results), just via a different,
% independent method, not this shared system. Both facts are true at once, not a contradiction.
\begin{table}[H]
\centering
\textbf{Feature-set purpose by question}\\[4pt]
\scriptsize
\begin{tabularx}{\linewidth}{@{} p{2.6cm} >{\hsize=0.75\hsize\raggedright\arraybackslash}X p{2.1cm} p{1.5cm} >{\hsize=1.25\hsize\raggedright\arraybackslash}X @{}}
\toprule
Feature set & Included feature groups & Model comparison & Environmental attribution & Reason \\
\midrule
Q3 (RSQ1) Set1--Set5 & Baseline + terrain/wind/climate/soil/edge (VIF-screened) & Yes (primary) & No & Feeds \texttt{dnn\_env\_terrain}/\texttt{pinn\_env\_terrain*}, designed for a fair DNN-vs-PINN comparison under progressively richer environmental input, and to test whether environmental conditioning improves height prediction. No linear/coefficient model reads these; not designed for attribution. \\
Q1 (RSQ2) Set1--Set5 & Baseline + terrain/wind/climate/soil/edge (VIF-screened) & Secondary (Elastic Net vs.\ XGBoost vs.\ LMM/NLME) & Yes (primary) & Feeds Elastic Net/XGBoost/LMM on \texttt{mean\_cr\_residual}, designed to identify which conditions align with persistent departure from the shared CR curve. Model comparison across the three fitted models is possible but secondary to the attribution question. \\
Q2 (RSQ3) Set1--Set5 & Baseline + terrain\_wind/climate/soil\_site/edge\_position (VIF-screened, one-hot, reference level dropped) & Secondary (Elastic Net vs.\ XGBoost vs.\ GNNWR) & Yes (primary) & Feeds Elastic Net/XGBoost/GNNWR on \texttt{local\_y\_max\_difference}, per-plot attribution; GNNWR specifically tests whether attribution is spatially varying, not just globally linear. \\
Q2 (RSQ3) \texttt{SCOPE\_GROUPS} (pre-existing, untouched) & Named category combinations (\texttt{terrain\_wind}, \texttt{+climate}, \texttt{+soil\_site}, \texttt{+edge\_position}, \texttt{+management}, \texttt{broad\_environment}, \texttt{broad\_environment\_plus\_management}) & No & Yes (category-level) & A \emph{different}, coarser question, "does adding this whole category help", answered by a pre-existing system this work did not modify. No Q1 equivalent of this specific system exists, though Q1 answers the same coarser question via its own separate category-ablation method (see Q1 appendix). \\
6-survey informational check (Q1/RSQ2 Set2 only) & Same candidate pool, 6-survey cohort & No & Diagnostic only & Not a primary result, checks whether Set2's top-10 changes on the smaller (7,467-row), noisier 6-survey cohort. It does (only 1 of 10 candidates overlap with 4-survey's Set2), reported as a caveat on Set2's stability, not as a second deliverable. \\
\bottomrule
\end{tabularx}
\caption{\small{What each feature set is for, and which question(s) it serves. Widening a Set (Set2$\to$Set3$\to$Set4$\to$Set5) is not claimed to necessarily improve prediction or strengthen an attribution finding, each Set is a candidate for downstream evaluation, and which one performs best is an empirical question for the project's real five-fold spatial CV, not decided by this screening pipeline.}}
\label{tab:feature-set-purpose}
\end{table}


\chapter{Appendix: Results}

\paragraph*{Q1 appendix}
\label{app:results-q1}

\paragraph*{Full Set4 variable table (EN vs.\ LMM, static/dynamic, stability)}
\label{app:nonlinearity-check}
% Referenced 3x from results_q1_21_08_0754pm.tex (the "one mechanism" paragraphs).
% Source: outputs/spatial_block_kfold/rq2_attribution_nested_set4_gated_all_vif/4survey/fold_{0-4}/,
% built 2026-08-21, see that file's own header comment for full provenance.

{\scriptsize
\setlength{\tabcolsep}{2pt}
\renewcommand{\arraystretch}{1.05}
\begin{longtable}{p{2.6cm}p{1.5cm}p{1.5cm}p{1.3cm}p{1.7cm}p{1.6cm}p{2.6cm}}
\caption{\small{Set4 (4-survey), all 19 variables: Elastic Net vs.\ mixed-effects (LMM) coefficients (mean $\pm$ SD, 5 folds), whether they agree in sign, static/dynamic classification, and stability verdict. ``Same sign?'' is a simplified same-fold-set check, not identical to the fuller stable/unstable definition used in-text (which also considers cross-set consistency and coefficient-vs-SD magnitude), topex and soilgrids\_ph pass this simplified sign check in Set4 alone but are still classed unstable by the fuller definition. Rows marked ``Not assessed'' were never individually checked for stability in the main text; their EN/LMM sign agreement here is offered as a first-pass screen, not a final verdict.}}
\label{tab:q1-set4-full-variable-table} \\
\toprule
Variable & EN coef.\ (mean$\pm$SD) & LMM coef.\ (mean$\pm$SD) & Same sign? & Static/Dynamic & Stability verdict & Basis \\
\midrule
\endfirsthead
\multicolumn{7}{l}{\normalfont\itshape (continued from previous page)} \\
\toprule
Variable & EN coef.\ (mean$\pm$SD) & LMM coef.\ (mean$\pm$SD) & Same sign? & Static/Dynamic & Stability verdict & Basis \\
\midrule
\endhead
\bottomrule
\multicolumn{7}{r}{\normalfont\itshape (continued on next page)} \\
\endfoot
\bottomrule
\endlastfoot
CanopyCover & $+1.943\pm0.038$ & $+1.864\pm0.078$ & Yes & Dynamic & Not assessed (leakage check instead) &, \\
time\_since\_thinning & $-2.139\pm0.216$ & $-1.475\pm0.596$ & Yes & Dynamic & Not assessed &, \\
time\_since\_thinning\_missing & $-2.995\pm0.182$ & $-1.321\pm0.698$ & Yes & Dynamic & Not assessed &, \\
recent\_thinning\_5yr & $-0.975\pm0.273$ & $-0.403\pm0.286$ & Yes & Dynamic & Not assessed &, \\
slope\_degrees & $+0.993\pm0.217$ & $+0.861\pm0.238$ & Yes & Static & \textbf{Stable} & Text-discussed \\
topex & $+0.063\pm0.296$ & $+1.325\pm0.155$ & Yes (Set4 only; flips Set3$\to$Set4) & Static & \textbf{Unstable} & Text-discussed; real local signal (0.09$\to$0.21) \\
tas\_mean & $-0.334\pm0.244$ & $+0.016\pm0.358$ & \textbf{No} & Dynamic & \textbf{Unstable} & Text-discussed; artifact (0.27$\to$0.08) \\
soilgrids\_ph & $+0.123\pm0.186$ & $+0.068\pm0.225$ & Yes (Set4 only) & Static & \textbf{Unstable} & Text-discussed; artifact (0.27$\to$0.05) \\
chelsa\_bio12\_precip\_mm & $-0.620\pm0.505$ & $-1.417\pm0.618$ & Yes & Static & \textbf{Stable} (EN/LMM agree) & Text-discussed; local-signal retention ($-0.298\to-0.0996$, 33\%) is close to the confirmed-artifact tas\_mean's (32\%), stability here rests on EN/LMM agreement itself, not a clearly demonstrated surviving local signal like topex's (267\%) \\
ceh\_twi & $-0.038\pm0.174$ & $+0.053\pm0.047$ & \textbf{No} & Static & \textbf{Unstable (new)} & Found this session; artifact ($-0.10\to0.01$) \\
gwa\_weibull\_a\_50m & $-0.579\pm0.609$ & $+0.342\pm0.066$ & \textbf{No} & Static & \textbf{Unstable (new)} & Found this session; real local signal ($-0.34\to-0.14$) \\
local\_relief\_500m & $+0.163\pm0.110$ & $-0.122\pm0.254$ & \textbf{No} & Static & \textbf{Unstable (new)} & Found this session; artifact ($-0.15\to0.00$) \\
gwa\_weibull\_k\_50m & $-0.135\pm0.196$ & $-0.359\pm0.249$ & Yes & Static & Not assessed &, \\
gwa\_weibull\_a\_10m & $-0.082\pm0.094$ & $-0.111\pm0.056$ & Yes & Static & Not assessed &, \\
gwa\_wind\_speed\_10m & $-0.121\pm0.074$ & $-0.108\pm0.068$ & Yes & Static & Not assessed &, \\
eastness & $+0.471\pm0.087$ & $+0.326\pm0.019$ & Yes & Static & Not assessed &, \\
dist\_to\_scpt\_boundary & $-0.299\pm0.103$ & $-0.146\pm0.071$ & Yes & Static & Not assessed &, \\
dist\_to\_road & $-1.332\pm0.171$ & $-1.216\pm0.130$ & Yes & Static (vintage uncertain) & Not assessed (own role in leakage check) &, \\
whcl & $-0.594\pm0.228$ & $-0.459\pm0.400$ & Yes & Static (not verified) & Not assessed &, \\
\end{longtable}
}

\begin{figure}[H]
\centering
\includegraphics[width=0.6\textwidth]{figures/fig_appendix/q1_final_appendix_vif.png}
\caption{\small{Variance inflation factor (VIF) for all 19 Set4 columns (4-survey), computed live via \texttt{compute\_vif\_table()}. Only the 3 thinning/management columns exceed the standard VIF=5 cutoff (expected, \texttt{time\_since\_thinning} and its missingness flag are closely related by construction); all abiotic (terrain/climate/soil) variables discussed in the main text, including the 6 unstable ones, have VIF well under 5, ruling out multicollinearity as their cause. \texttt{CanopyCover} itself also has low VIF despite being the strongest predictor of the target: VIF measures how well a variable is explained by the \emph{other} predictors, not how strongly it predicts the target, and CanopyCover (a stand-structure attribute) is not well explained by the terrain, wind, climate, and soil columns around it. This is the same conclusion as the earlier non-redundancy check (Figure~\ref{fig:q1-fig1-canopycover-ablation}) restated in different terms: CanopyCover carries information the other 18 columns cannot substitute for.}}
\label{fig:q1-appendix-vif}
\end{figure}

\paragraph*{Category-level ablation (soil / climate / terrain)}
\label{app:category-ablation}
% AUDIT: still a placeholder. A candidate source (outputs/spatial_block/rq3_category_attribution/
% 4survey/category_morans_i_before_after.csv) was checked 2026-08-25 and ruled out: it gives soil
% +0.025 R2 (close, plausibly the same result) but climate +0.021 R2, i.e. climate HELPS in that
% run, which contradicts the main text's "-0.056, actively reduces accuracy" claim directly, not
% just imprecisely. Different category definitions or feature scope, not the same analysis under a
% legacy name. Do not cite that CSV for this claim without first reconciling why the sign flips.
Removing soil variables (\texttt{soilgrids\_ph}, \texttt{ceh\_pedotope},
\texttt{ceh\_subsurface\_drainage}, \texttt{ceh\_textural\_composition},
\texttt{dist\_to\_watercourse}) from Set4 changes test $R^2$ by $+0.022$; removing macroclimate
variables (\texttt{tas\_mean}, \texttt{groundfrost\_mean}, \texttt{chelsa\_bio1\_celsius},
\texttt{chelsa\_gdd5\_degc}, \texttt{chelsa\_bio12\_precip\_mm}) changes it by $-0.056$, i.e.\
climate variables reduce held-out accuracy when present. These two headline deltas are the
figures already quoted in the main Q1 Discussion. The full refit table behind them (Set4 vs.\
each category individually removed, all five categories) has not been re-verified against a
saved source file this session; a fuller breakdown table will be added here once that source
run is reproduced or reconciled with the mismatched candidate noted above.

\paragraph*{Q2 appendix}
\label{app:results-q2}

% Referenced from results_q2_22_08_0214.tex (the "same terrain variables behave the same
% way" cross-check paragraph). Built 2026-08-23.

\begin{table}[H]
\centering
\textbf{Terrain/climate variable reliability, Q1 vs.\ Q2}\label{app:results-q2-reliability}\\[4pt]
\scriptsize
\setlength{\tabcolsep}{3pt}
\renewcommand{\arraystretch}{1.2}
\begin{tabular}{p{2.4cm}p{2.6cm}p{3.6cm}p{1.3cm}p{2.6cm}}
\toprule
Variable & Q1: EN vs.\ LMM & Q1 mechanism (if disputed) & Q2: \% cpmt.\ same sign & Read \\
\midrule
CanopyCover & Agree (EN 1.94, LMM 1.86) & not disputed & 100.0 & Reliable everywhere \\
slope\_degrees & Agree (EN 0.99, LMM 0.86) & not disputed & 97.0 & Reliable everywhere \\
chelsa\_bio12\_precip\_mm & Agree (EN $-0.62$, LMM $-1.42$) & not disputed & 87.4 & Reliable everywhere \\
topex & Disagree (EN 0.06, LMM 1.32) & Real local signal (within-cpmt.\ corr.\ 0.214 $>$ raw 0.080) & 71.9 & Real effect, hard for a straight line (global or local) to pin down consistently \\
windward\_topex & not tested in Q1 & n/a & 91.8 & Q2-only variable, more stable than plain topex \\
tas\_mean & Disagree (EN $-0.33$, LMM 0.02) & Artifact (within-cpmt.\ corr.\ 0.090 $\ll$ raw 0.279) & 78.4 & Q1: fake signal, despite a fairly stable local sign in Q2, stability alone does not make it a real effect \\
soilgrids\_ph & Disagree (EN 0.12, LMM 0.07) & Artifact (within-cpmt.\ corr.\ 0.035 $\ll$ raw 0.224) & 51.1 & Both chapters agree: not a real effect \\
dist\_to\_road & not tested in Q1 & n/a & 74.9 & No Q1 comparison; moderately stable sign in Q2 \\
\bottomrule
\end{tabular}
\caption{\small{Terrain/climate variable reliability, cross-referenced between Q1 (does a
global Elastic Net model agree with a linear mixed-effects model on a variable's
coefficient, Set4, 4-survey) and Q2 (does GNNWR's own local coefficient keep the same sign
across compartments, Set4, 4-survey, pooled 5 folds, $n=231$ compartments). These are two
genuinely different tests that happen to agree on topex and slope\_degrees, not the same
measurement repeated.}}
\label{tab:results-q2-reliability}
\end{table}
% Source (internal tracking only, not for the caption): Q1 columns live-recomputed from
% outputs/spatial_block_kfold/rq2_attribution_nested_set4_gated_all_vif/4survey/fold_{0-4}/
% (elastic_net_coefficients.csv, nlme_fixed_effects.json); Q2 column live-recomputed from
% GNNWR's own saved per-plot coef_* columns, same fold files as Table~\ref{tab:results-q2-ablation}.

% Source: notebooks/results_q2/results_figures_q2_example_plot.ipynb, cells 9-11. Plot 166972
% chosen as the closest-to-median-|residual| candidate among non-flagged plots with >=4
% surveys (the same presentability requirement as the main-text example), not the best-fit
% tail, the point of this figure is to show what attribution looks like for an ordinary
% plot, not another flattering case.

\begin{figure}[H]
\centering
\textbf{A typical example plot}\label{app:results-q2-typical-plot}\\[4pt]
\includegraphics[width=\textwidth]{figures/fig_appendix/q2_example_plot_curve_typical_appendix.png}
\caption{\small{Plot 166972, chosen because its prediction error sits at the population
median, not the well-fit tail (ages 35--50, heights 26.0--37.3\,m). Left: observed trajectory
against its own fitted curve and the yield-class benchmark curve. Right: XGBoost SHAP
contributions to the predicted $\Delta y_{\max}$. The actual deviation is $+11.2$\,m but
XGBoost predicts only $+5.3$\,m, roughly half, with a diffuse SHAP breakdown
($\mathit{chelsa\_bio12\_precip\_mm}$, CanopyCover, $\mathit{recent\_thinning\_5yr}$,
$\mathit{dist\_to\_road}$, $\mathit{windward\_topex}$, none dominant). Companion to
Figure~\ref{fig:q2-example-plot}'s Plot 309770: that plot was chosen because it is unusually
well-explained, and this one shows what the same attribution mechanism looks like for an
ordinary plot instead.}}
\label{fig:q2-example-plot-typical}
\end{figure}

\paragraph*{Flagged-plot split diagnostics}
\label{app:flagged-plot-split}
% Source: q2_brain_dump_21_08.md, sections 4-5. Not a new analysis, transcribing already-run
% diagnostics into the appendix, including the open items they themselves disclosed.
Around 10 plots recur as the worst residuals across every model and feature set (a flag rate
roughly 36$\times$ higher than chance, not a random tail). This traces to the fixed-$k$
curve-fitting mechanism used to build $\Delta y_{\max}$ (Equation~\ref{eq:local-ymax-diff}): growth rate is
pinned to the plot's assigned yield class, so a plot that genuinely grows slower or less
smoothly than that assumption forces an inflated ceiling estimate instead. Scaling this up with
Tukey's $1.5\times\text{IQR}$ fence flags 266 plots in the 4-survey cohort (0.47\%) and 709 in
the 6-survey cohort (5.32\%), clustered by compartment. Removing these plots from training does
not improve $R^2$, they carry real signal, not just noise.

A population-wide split test (recomputed on 250 of these plots) asked whether each flagged
plot's pattern looks like the fixed-$k$ artefact or like genuinely fast growth: 71.6\% fit the
artefact signature (exposed terrain), and 28.4\% look like genuinely fast, real growth
(sheltered terrain), tested with actual $p$-values rather than judged by eye. Two open items are
disclosed rather than resolved here: the 266-vs-250 count difference between the Tukey-fence
count and the split-test sample has not been traced to the exact plots involved, though both
counts were independently verified; and the artefact mechanism is confounded with a real
wind-exposure effect, since exposed terrain plausibly drives both genuinely unusual growth and
the artefact pattern, and this data cannot fully separate the two. This is the same
rate-ceiling-mixing mechanism described for the whole population in
Section~\ref{sec:eval-robustness}, just at a far more extreme scale for this flagged tail,
not a second, unrelated finding.

\paragraph*{Category-level Moran's I diagnostics (with terrain bandwidth-artifact adjustment)}
\label{app:results-q2-category-morans-i}
% AUDIT: placeholder, results_q2 currently points this content at the generic \ref{chap:appendix} rather than
% a specific label; once this subsection is filled in, update that in-text ref to
% \ref{app:results-q2-category-morans-i} instead.

\paragraph*{Q3 appendix}
\label{app:results-q3}

\paragraph*{XGBoost hyperparameter search}
\label{app:hyperparameters}
% Source: TEMP_results/TEMP_rq1_xgb_hyperparameter_search_2026-08-16.tex. Despite the "rq1"
% filename, this search retunes RQ1's XGBoost script but on RQ3's OWN copy of the same fitting
% code and reports numbers that match Q3's headline Table (tab:results-q3) row exactly, so this
% is genuinely Q3's search, not a mislabelled Q1 result, verified against the winning-config
% test R2/RMSE/MAE, which match tab:results-q3's XGBoost (tuned) row to 3 decimal places.
Grid: \texttt{n\_estimators} $\in \{300, 500, 800\}$, \texttt{max\_depth} $\in \{3, 4, 6\}$,
\texttt{learning\_rate} $\in \{0.02, 0.04, 0.08\}$, 27 configurations, each fit on Set3 with
5-fold spatial-block cross-validation (135 fits per cohort, 270 total across both cohorts).
Model selection used mean validation $R^2$ only, never test $R^2$, to avoid tuning-on-test
leakage; test metrics are reported only for the winning configuration, after selection was
already decided.

\begin{table}[H]
\centering
\scriptsize
\renewcommand{\arraystretch}{1.1}
\begin{tabular}{lrrrrrrr}
\toprule
Cohort & $n_{\mathrm{est}}$ & Depth & LR & Val.\ $R^2$ & Test $R^2$ & Test RMSE & Test MAE \\
\midrule
4-survey (winning) & 500 & 6 & 0.02 & 0.682 & $0.674\pm0.009$ & $4.548\pm0.137$ & $3.431\pm0.123$ \\
4-survey (borrowed default) & 500 & 4 & 0.04 & 0.678 & $0.675\pm0.009$ & $4.539\pm0.152$ & $3.431\pm0.153$ \\
6-survey (winning) & 300 & 3 & 0.08 & 0.717 & $0.722\pm0.031$ & $3.656\pm0.232$ & $2.669\pm0.183$ \\
6-survey (borrowed default) & 500 & 4 & 0.04 & 0.712 & $0.718\pm0.025$ & $3.682\pm0.221$ & $2.697\pm0.170$ \\
\bottomrule
\end{tabular}
\caption{\small{Winning configuration per cohort (selected on mean validation $R^2$ across 5
folds) against the generic default config borrowed from the same fitting code before this
search was run. The 4-survey winning-by-validation config is marginally \emph{lower} on test
$R^2$ than the borrowed default (0.674 vs.\ 0.675), expected noise from selecting on
validation rather than test, not a sign the search failed. Across both cohorts, the full
270-fit search moves test $R^2$ by at most 0.004, well inside the fold-to-fold SD (0.009 on
4-survey, 0.025--0.031 on 6-survey): the borrowed default was already close to optimal, and
XGBoost's accuracy edge over the neural models is not an artefact of a lucky or unfair
hyperparameter choice.}}
\label{tab:xgb-hyperparameter-search}
\end{table}

% Moved out of the main Table~\ref{tab:results-q3} 2026-08-23 for clarity, since 6-survey is
% not this chapter's main analysis surface (see that table's caption). PINN/PINN-k rows show
% --: the corrected-forward-pass work this chapter reports (environment sweep, lambda
% ablation, hyperparameter checks, see temp_results_pinn/RESULTS_TABLE.md) was all done on
% 4-survey only; 6-survey was never rerun with the corrected architecture, and the pre-fix
% 6-survey numbers are not valid evidence for the corrected model (same reasoning as the
% pre-fix PINN runs excluded from the main table).

\begin{table}[H]
\centering
\textbf{Full 6-survey table}\label{app:results-q3-6survey}\\[4pt]
\scriptsize
\setlength{\tabcolsep}{4pt}
\renewcommand{\arraystretch}{1.2}
\begin{tabular}{lrrr}
\toprule
Model & RMSE${\downarrow}$ & MAE${\downarrow}$ & $R^2{\uparrow}$ \\
\midrule
Linear & $4.877{\pm}0.643$ & $3.619{\pm}0.426$ & $0.509{\pm}0.067$ \\
RF & $4.062{\pm}0.249$ & $3.010{\pm}0.185$ & $0.656{\pm}0.046$ \\
XGBoost (tuned) & $\mathbf{3.656{\pm}0.232}$ & $\mathbf{2.669{\pm}0.183}$ & $\mathbf{0.722{\pm}0.031}$ \\
DNN & $3.903{\pm}0.316$ & $2.867{\pm}0.275$ & $0.684{\pm}0.031$ \\
PINN &, &, &, \\
PINN-$k$ &, &, &, \\
\bottomrule
\end{tabular}
\caption{\small{Model comparison for top-height prediction, 6-survey, Set3,
spatial\_block\_kfold, per-fold mean $\pm$ SD across five folds. n=82{,}614 row-years
(13{,}769 plots). Referenced from the main text's DNN-vs.-XGBoost fold-count and
compartment-subsampling discussion. PINN/PINN-$k$ were not rerun with the corrected
forward pass on this cohort (--); see main-text Table~\ref{tab:results-q3}'s caption.}}
\label{tab:results-q3-6survey}
\end{table}

\label{app:results-q3-tables}
% Moved out of the main text 2026-08-24, main Q3 results section is capped at 3 tables
% (Table~\ref{tab:results-q3}, Table~\ref{tab:pinn-conditions}, Table~\ref{tab:plausibility-comparison}
% stay in main text). These six are supporting/diagnostic detail, referenced from the main-text
% discussion via \ref{} as normal, moving them here does not change any in-text reference.

% --- Table: compartment 1129 cross-fold check (unchanged) ---
\begin{table}[H]
\centering
\textbf{Q3 results tables}\\[4pt]
\scriptsize
\begin{tabular}{lrrr}
\toprule
Fold & Split membership & Mean $y_{max}$ deviation & Rows $>$70\,m (of 1484) \\
\midrule
0 & 100\% test & $+12.1$\,m & 72 \\
1 & 95\% train & $+11.5$\,m & 44 \\
2 & 98\% train & $+10.7$\,m & 0 \\
3 & 100\% train & $+15.6$\,m & 0 \\
4 & 100\% validation & $+12.4$\,m & 0 \\
\bottomrule
\end{tabular}
\caption{\small{Compartment 1129's mean $y_{max}$ deviation in each of the 5 spatial-block
folds, by how much of the compartment's own data that fold trained on. Inflation persists (and
does not shrink) even when the model trains on the compartment's own real data.}}
% source: RESULTS_TABLE.md #10
\label{tab:compartment1129-crossfold}
\end{table}

% --- Table: trunk-residual comparison, plain PINN vs. PINN-k ---
\begin{table}[H]
\centering
\scriptsize
\begin{tabular}{lrrr}
\toprule
Fold & Plain PINN trunk mean$|\cdot|$ & PINN-$k$ trunk mean$|\cdot|$ & Ratio \\
\midrule
0 & 0.207 & 0.410 & 1.98$\times$ \\
1 & 0.238 & 0.553 & 2.33$\times$ \\
2 & 0.255 & 0.844 & 3.31$\times$ \\
3 & 0.339 & 0.668 & 1.97$\times$ \\
4 & 0.268 & 0.704 & 2.63$\times$ \\
\bottomrule
\end{tabular}
\caption{\small{Mean absolute size of the shared correction network's output, plain PINN vs.\
PINN-$k$, all 5 folds. PINN-$k$'s correction is larger than plain PINN's in every fold,
roughly $2\times$ or more throughout, not a strict trend.}}
% source: RESULTS_TABLE.md #7. UPDATED 2026-08-24: added fold 3 (ratio 1.97x), completing all
% 5 folds. Fold 3 is just under 2x (1.968), so "consistently above 2x" softened to "roughly 2x
% or more", range across all 5 folds is 1.97x-3.31x, no fold below ~1.97x. Earlier "grows
% fold to fold" framing already dropped when fold 4 was added (1.98 -> 2.33 -> 3.31 looked
% monotonic on 3 folds; fold 4 at 2.63 broke it), fold 3 confirms this further: not a trend,
% just consistently >=~2x across all 5 folds. All 5 folds now done, nothing pending.
\label{tab:trunk-residual}
\end{table}

% --- Table: error by age band and height band (unchanged) ---
\begin{table}[H]
\centering
\small
\begin{tabular}{lrrrr}
\toprule
Age band (years) & $n$ & DNN MAE & PINN MAE & PINN-$k$ MAE \\
\midrule
15--30 & 29,837 & 2.49 & 2.52 & 2.53 \\
30--45 & 17,861 & 2.69 & 2.85 & 2.93 \\
45--60 & 7,590 & 4.16 & 4.24 & 4.42 \\
60--75 & 2,034 & 4.63 & 5.23 & 5.15 \\
75--93 & 751 & \textbf{6.83} & 6.40 & 6.81 \\
\midrule
Height band & $n$ & DNN MAE & PINN MAE & PINN-$k$ MAE \\
\midrule
0--10\,m & 12,373 & 2.68 & 3.37 & 3.15 \\
10--20\,m & 29,629 & 2.59 & 2.62 & 2.74 \\
20--30\,m & 13,733 & 3.36 & 2.98 & 3.15 \\
30--47\,m & 2,338 & 5.26 & 5.71 & 5.72 \\
\bottomrule
\end{tabular}
\caption{\small{Mean absolute error by tree age and by tree height, DNN vs.\ plain PINN vs.\
PINN-$k$, pooled across all 5 folds (58,073 plots for all three models). Bold: DNN is the
WORST of the three at the oldest ages, not the best, see note below on how this reverses the
fold-0-only reading.}} % source: RESULTS_TABLE.md #12, updated 2026-08-24, pooled 5-fold
\label{tab:age-height-error}
\end{table}

% --- Table (NEW): selection rule + the two flagship examples ---
% Rule (2026-08-23, replaces earlier ad-hoc picks): among compartments with >=50 pooled
% held-out plots (excludes small-n compartments where 1-2 plots could drive the mean), rank by
% mean deviation for the "systematic site signal" story, and separately by single-plot max/min
% deviation for the "one implausible plot" story. A compartment can lead on one ranking and not
% the other, that is expected, not a problem, since the two rankings answer different
% questions.
\begin{table}[H]
\centering
\small
\begin{tabular}{llrrp{3.3cm}}
\toprule
Compartment & Direction & Mean dev. (m) & Most extreme plot (m) & Why flagged \\
\midrule
1129 & Over & +12.2 (rank 10/178) & \textbf{+25.8 (rank 1/178)} & Single most extreme plot \\
2021 & Under & \textbf{-21.4 (rank 1/178)} & \textbf{-30.4 (rank 1/178)} & Largest mean \emph{and} most extreme plot \\
\bottomrule
\end{tabular}
\caption{\small{The two flagship compartments, ranked among the 178 compartments with
$n\geq50$ pooled held-out plots. 1129 wins on single-plot extremity alone; 2021 wins on both
mean and single-plot extremity.}}
% source: RESULTS_TABLE.md #13, n>=50 reliability rule applied to the pooled 5-fold compartment
% deviation table (58,073 plots).
\label{tab:productivity-flagships}
\end{table}

% --- Table (NEW): the wider top-5-per-side picture, reported honestly (not filtered to
% only the yield-class-confirming rows) ---
\begin{table}[H]
\centering
\small
\begin{tabular}{lrrrlr}
\toprule
Compartment & Mean dev. (m) & $n$ plots & Yield class & Agrees? & Mean age (yr) \\
\midrule
\multicolumn{6}{l}{\textit{Over-productive, top 5 by mean deviation (n$\geq$50)}} \\
2057 & +15.6 & 112 & 10--16 (mixed) & No & 59.8 \\
1130 & +14.6 & 78 & 22--24 & Yes & 33.6 \\
2130 & +14.6 & 98 & 20--24 & Yes & 35.1 \\
1122 & +13.4 & 51 & 24 & Yes & 25.0 \\
1027 & +13.4 & 56 & 12 & No & 72.0 \\
\midrule
\multicolumn{6}{l}{\textit{Under-productive, top 5 by mean deviation (n$\geq$50)}} \\
2021 & -21.4 & 403 & 12 & Yes &, \\
1070 & -18.1 & 222 & 12 & Yes &, \\
2094 & -17.7 & 279 & 10--12 & Yes &, \\
1069 & -11.8 & 297 & 12 & Yes &, \\
2022 & -10.3 & 417 & 16 & Weakly &, \\
\bottomrule
\end{tabular}
\caption{\small{Top 5 over- and under-productive compartments by mean deviation
($n\geq50$), checked against independent yield class. Reported as-is, not filtered down to the
rows that agree. Under-productive side agrees with independent yield class in 4/5 cases;
over-productive side agrees in only 3/5, including a disagreement from the single biggest
over-productive signal (2057). The two disagreeing over-productive rows (2057, 1027) are both
old and wind-exposed (mean age 59.8/72.0yr, windward index +15.2/+11.0); the three agreeing
rows (1130, 2130, 1122) are both young and sheltered-to-neutral (mean age 25.0--35.1yr,
windward index -9.5 to +0.4), see main-text discussion of this split.}}
% source: RESULTS_TABLE.md #13.
\label{tab:productivity-top5}
\end{table}

% --- Table (NEW): correlation with independent yield-class records ---
\begin{table}[H]
\centering
\small
\begin{tabular}{lrr}
\toprule
Comparison & Pearson $r$ & $n$ \\
\midrule
Plain PINN $y_{max}$ deviation vs.\ yield class & 0.273 & 58,073 \\
PINN-$k$'s $k$ vs.\ yield class & 0.286 & 58,073 \\
PINN-$k$'s $y_{max}$ vs.\ yield class & 0.135 & 58,073 \\
\bottomrule
\end{tabular}
\caption{\small{Correlation between each model's personalisation and the Forestry Commission's
independent yield-class rating, never a model input, checked against both feature lists.
Pooled across all 5 folds.}}
% UPDATED 2026-08-24: was fold-0-only (n=11,508, PINN-k y_max r=0.334, the apparent
% strongest of the three). Pooled data flips the ranking: PINN-k's y_max is now the WEAKEST
% (0.135), its k the strongest (0.286). See RESULTS_TABLE.md section 13 addendum.
\label{tab:yieldclass-correlation}
\end{table}

\label{app:results-q3-figures}
% Moved out of the main text 2026-08-24, same reasoning as Section~\ref{app:results-q3-tables}:
% main Q3 Discussion keeps only the two headline compartment maps (fig:ymax-deviation-map,
% fig:ymax-k-deviation-maps) in the flow; these three are reference/ranking detail one step
% removed from the two flagship examples, not themselves the headline claim. Referenced from the
% main text via \ref{} as normal, moving them here does not change any in-text reference.

\begin{figure}[H]
\centering
\textbf{Q3 supporting figures}\\[4pt]
\includegraphics[width=0.95\textwidth]{figures/fig_appendix/q3_pinn_param_distribution.png}
\caption{\small{PINN-$k$'s per-plot parameters, pooled across all 5 folds (58,073 plots), each
plot's deviation measured against its own fold's population baseline. Left: $y_{max,i}$ stays
tightly pinned near the population value for the large majority of plots (76.7\% within
$\pm$1m). Right: $k_i$ is skewed strongly toward faster than population for most plots (92.1\%
moderately or much faster), a population-wide shift, not a symmetric spread of real
heterogeneity. Quantitative companion to Figure~\ref{fig:ymax-k-scatter}'s shape: the same
one-sided $k$ routing, shown as a distribution rather than a scatter.}}
% UPDATED 2026-08-24: pooled across all 5 folds (was fold-0-only, n=11,508), same fold-own-
% baseline pooling as load_pooled_pinn_k_deviation_data() in build_q3_redraft_figures.py, to
% avoid the same Simpson's-paradox-style confound already found and fixed in the y_max-vs-k
% scatter figure. Moved into figures/fig_results/q3/ at the same time (was in the flat
% fig_results/ root, unreferenced anywhere). Script:
% models/growth_curve_attribution/q3_pinn_param_distribution_check.py.
\label{fig:pinn-param-distribution}
\end{figure}

\begin{figure}[H]
\centering
\includegraphics[width=0.95\textwidth]{figures/fig_appendix/q3_all_compartments_reference_map.png}
\caption{\small{Every compartment named anywhere in this chapter, numbered on one map, with
the reason each was highlighted. A single lookup source, some compartments here appear in a
table but no other figure, or a figure but no table.}}
% BUILT 2026-08-24. Reference/lookup figure, not tied to one specific paragraph.
\label{fig:all-compartments-reference}
\end{figure}

\begin{figure}[H]
\centering
\includegraphics[width=0.95\textwidth]{figures/fig_appendix/q3_pinn_ymax_hotspot_summary.png}
\caption{\small{Top-3 over- and under-productive compartments (n$\geq$50), plain PINN (left)
vs.\ PINN-$k$ (right). 2057 tops both; the under-productive top-3 is the same three compartments
on both sides, just reordered.}}
\label{fig:hotspot-summary}
\end{figure}

\begin{figure}[H]
\centering
\includegraphics[width=0.7\textwidth]{figures/fig_appendix/q3_yieldclass_scatter.png}
\caption{\small{Plain PINN's $y_{max}$ deviation vs.\ independent yield class, pooled across
all 5 folds. Green: every held-out test plot (58,073), one dot each, yield class only takes
about a dozen discrete values, which is why the dots stack into vertical bands rather than
spreading evenly across the x-axis. Red and blue: the same kind of point, restricted to the two
flagship compartments (1129, 2021), so their position within the wider pattern is visible,
1129 clustered high (mostly yield class 22-24), 2021 clustered low and tightly (entirely yield
class 12, strongly negative). The dashed black line is a zero-deviation reference, not a data
series. Real but modest correlation (Table~\ref{tab:yieldclass-correlation}), the upward
drift of the whole cloud left to right is what that correlation number is measuring.}}
% REBUILT 2026-08-24 on pooled data (was fold-0-only, n=11,508). Pooling changed the answer
% here more than expected, see RESULTS_TABLE.md and the note on fig:ymax-k-scatter's build
% comment for why pooling is the correct choice for THIS comparison (yield class has no
% fold-varying baseline to confound, unlike the y_max-vs-k comparison). Source: chat 2026-08-24.
\label{fig:yieldclass-scatter}
\end{figure}

\paragraph*{Why the broadest feature list looked worse}
\label{app:results-q3-set4}
% Moved out of the main Discussion 2026-08-24, secondary caveat, not a headline. Was its own
% \paragraph* in the working file, titled "Appendix: why the broadest feature list looked worse
%, and why that is a weak finding, not a strong one."
Why the broadest environment list (Set4) seemed to underperform the medium list (Set3): its
extra features are not ignored by the model, they draw more than their fair share of
permutation importance, they just do not repay that attention in held-out accuracy as well as
the medium list's more curated features. Caveat: the Set3-vs-Set4 gap this explains is itself
within fold-to-fold noise for plain PINN (Table~\ref{tab:pinn-conditions}), and absent entirely
for PINN-$k$, a plausible explanation for a small, uncertain effect, not a solid one.

\paragraph*{Pre-fix PINN runs (uncorrected baseline, for comparison)}
\label{app:results-q3-prefix-pinn}
% AUDIT: placeholder, results_q3 currently points this content at the generic \ref{chap:appendix} rather than
% a specific label; once this subsection is filled in, update that in-text ref to
% \ref{app:results-q3-prefix-pinn} instead.


\begin{figure}[H]
\centering
\includegraphics[width=0.65\textwidth]{figures/fig_appendix/q3_pinn_ymax_k_scatter.png}
\caption{\small{Covariance between personalised parameters in PINN-$k$ (pooled across 5 folds). 
Plot-level asymptote departures ($\Delta y_{\max}$) plotted against growth-rate departures ($\Delta k$) relative to each fold's regional baseline ($r = 0.43$). The positive relationship demonstrates coordinated biological scaling across diverse stand conditions.}}
\label{fig:ymax-k-scatter}
\end{figure}
\end{document}
%=======================================================================================

